\documentclass[aspectratio=169, xcolor=dvipsnames]{beamer}

\AtBeginSection[]
{
    \begin{frame}
        \frametitle{Visão Geral}
        \tableofcontents[currentsection]
    \end{frame}
}

\usefonttheme{professionalfonts} % using non standard fonts for beamer
\usefonttheme{serif} % default family is serif
\usepackage{fontspec}
\setmainfont{XCharter}[
    UprightFont    = *-Roman,
    BoldFont       = *-Bold,
    ItalicFont     = *-Italic,
    BoldItalicFont = *-BoldItalic,
]

\usepackage{hyperref}
\usepackage{graphicx} % Allows including images
\usepackage{booktabs} % Allows the use of \toprule, \midrule and \bottomrule in tables
\input{./slides/SimplePlusTheme/beamerthemeSimplePlus.sty}
\input{./slides/SimplePlusTheme/beamerinnerthemeSimplePlus.sty}
\input{./slides/SimplePlusTheme/beamerfontthemeSimplePlus.sty}
\input{./slides/SimplePlusTheme/beamercolorthemeSimplePlus.sty}

\usepackage{listings}
\setbeamercovered{transparent}

\usepackage{biblatex}
\addbibresource{slides/reference.bib}
\renewcommand*{\bibfont}{\footnotesize}
% Tamanho dos slides de "Leituras Recomendadas": maior que a bibliografia
% completa (\bibfont acima, que rende \tiny via a redefinição de \footnotesize)
% para destacar as leituras curadas.
\newcommand{\leiturasize}{\small}

\usepackage[brazilian ]{babel}
\usepackage{csquotes}
\usepackage{pgfplots}
\pgfplotsset{compat=1.18}

\usepackage{emoji}
\usepackage{amsmath}
\usepackage[xcharter,bigdelims,vvarbb]{newtxmath}
% newtxmath's operators font requests the fontspec family `XCharter(0)' in OT1;
% without these substitutions math digits and operator names silently fall back
% to Computer Modern (LaTeX Font Warning: Font shape `OT1/XCharter(0)/m/n' undefined).
\DeclareFontFamily{OT1}{XCharter(0)}{}
\DeclareFontShape{OT1}{XCharter(0)}{m}{n}{<->ssub * XCharter-TLF/m/n}{}
\DeclareFontShape{OT1}{XCharter(0)}{m}{it}{<->ssub * XCharter-TLF/m/it}{}
\DeclareFontShape{OT1}{XCharter(0)}{b}{n}{<->ssub * XCharter-TLF/b/n}{}
\DeclareFontShape{OT1}{XCharter(0)}{bx}{n}{<->ssub * XCharter-TLF/b/n}{}

\usepackage{bm}
\usepackage[table]{xcolor}
\usepackage{xcolor}
\usepackage{algpseudocode}

\usepackage{cancel}
\usepackage{ulem}

\usetikzlibrary{shapes}
\usetikzlibrary{arrows}
\usetikzlibrary {arrows.meta}
\usetikzlibrary{calc,positioning}
\usepgfplotslibrary{fillbetween}
\usetikzlibrary{angles,quotes}
\usetikzlibrary{tikzmark}


\definecolor{PastelRed}{HTML}{FFC1CC}
\definecolor{PastelBlue}{HTML}{C2F0FF}
\definecolor{PastelBlue2}{HTML}{3182BD}

\definecolor{PastelPink}{HTML}{FFCBE1}
\definecolor{PastelGreen}{HTML}{D6E5BD}
\definecolor{PastelYellow}{HTML}{F9E1A8}
\definecolor{PastelBlue}{HTML}{BCD8EC}
\definecolor{PastelPurple}{HTML}{DCCCEC}
\definecolor{PastelOrange}{HTML}{FFDAB4}


\renewcommand{\footnotesize}{\tiny}

\newcommand{\mespor}{
  \ifcase\month
    \or janeiro
    \or fevereiro
    \or março
    \or abril
    \or maio
    \or junho
    \or julho
    \or agosto
    \or setembro
    \or outubro
    \or novembro
    \or dezembro
  \fi
}
% \usepackage{csquotes}
% \usepackage{minted}
% \usemintedstyle{xcode}
\definecolor{vscLightBackground}{HTML}{FFFFFF}
\definecolor{vscLightText}{HTML}{000000}
\definecolor{vscLightGreen}{HTML}{008000}        % For Comments AND Numbers
\definecolor{vscLightRed}{HTML}{A31515}          % For Strings
\definecolor{vscLightTeal}{HTML}{267F99}
\definecolor{vscBlue}{HTML}{0000FF}% For main keywords (if, for, def...)
\definecolor{vscLightPytorch}{HTML}{800080}      % For PyTorch keywords (dark magenta)

\lstdefinestyle{vscode-light}{
    language=Python,
    backgroundcolor=\color{vscLightBackground},
    basicstyle=\ttfamily\scriptsize\color{vscLightText},
    keywordstyle=\color{vscLightTeal},
    commentstyle=\color{vscLightGreen},
    stringstyle=\color{vscLightRed},
    % Style for PyTorch keywords (using keyword group 2)
    keywordstyle=[2]\color{vscLightPytorch}, 
    morekeywords={self, True, False, None},
    morekeywords=[2]{
        torch, Tensor, device, dtype, to, nn, optim, autograd, utils, data,
        Module, Linear, Conv2d, ReLU, CrossEntropyLoss, Adam, SGD,
        Dataset, DataLoader, randn, zeros, ones, arange, cat, stack,
        view, reshape, sum, mean, max, min, matmul, no_grad, backward
    },
    keywordstyle=[3]\color{vscBlue}, 
    morekeywords=[3]{
        model, loss, y, y_hat, x , optimizer
    },
    frame=single,
    tabsize=1,
    showstringspaces=false,
    breaklines=true,
    captionpos=b,
    extendedchars=true,
    numbers=left,
}

\usepackage[cache=true]{minted}
% minted v3+ (TeX Live 2024+) auto-detects the output directory.
% minted v2 (Ubuntu apt TeX Live 2023) needs it set explicitly to
% match latexmkrc's $out_dir = 'build'.
\makeatletter
\@ifpackagelater{minted}{2024/01/01}{}{%
  \def\minted@outputdir{build/}%
}
\makeatother
\setminted{
    fontsize=\scriptsize,
    style=vs,
    breaklines=true,
    numbers=left,
    numbersep=5pt,
    frame=single,
}
\providecommand{\citeauthoryear}[1]{(\citeauthor{#1},~\citeyear{#1})}

\providecommand{\thetav}{\bm{\theta}}

% vetor coluna 2D com colchetes
\providecommand{\cvec}[2]{\begin{bmatrix} #1 \\ #2 \end{bmatrix}}

\providecommand{\varvector}[1]{\mathbf{#1}}
% instance
\providecommand{\X}{\varvector{X}}
\providecommand{\mlinstance}{\varvector{X}}
\providecommand{\mlinstancei}{\varvector{X}^{(i)}}
\providecommand{\mlinstancex}[1]{\varvector{X}^{(#1)}}

% target
\providecommand{\y}{\varvector{y}}
\providecommand{\mltarget}{\varvector{y}}
\providecommand{\mltargeti}{\varvector{y}^{(i)}}
\providecommand{\mltargetx}[1]{\varvector{y}^{(#1)}}

%hats
\providecommand{\yhat}{\hat{\varvector{y}}}
\providecommand{\mltargethat}{\hat{\varvector{y}}}
\providecommand{\mltargethati}{\hat{\varvector{y}}^{(i)}}
\providecommand{\mltargethatx}[1]{\hat{\varvector{y}}^{(#1)}}

% linear reg
\providecommand{\mllinearregression}{\mltargethat=\mlinstance\bm{\theta}}

\providecommand{\tikzarrowcolor}[3]{%
  \begin{tikzpicture}[remember picture, overlay]
    % Arrow body
    \path[fill=#3] ([shift={(#1,#2)}]current page.south west) 
      rectangle ++(0.3, -0.2);
    % Arrow head
    \path[fill=#3] ([shift={(#1+0.3,#2+0.1)}]current page.south west) -- 
         ++(0, -0.4) -- 
         ++(0.2, 0.2) -- 
         cycle;
  \end{tikzpicture}%
}

\providecommand{\tikzarrow}[2]{%
    \tikzarrowcolor{#1}{#2}{Red}
}



\DeclareMathOperator*{\argmax}{argmax}
\DeclareMathOperator*{\argmin}{argmin}



%----------------------------------------------------------------------------------------
%    TITLE PAGE
%----------------------------------------------------------------------------------------

\title{Deep Learning}
\subtitle{\textit{Positional Encodings}}

\author{Lucas Silveira Kupssinskü}

\institute
{ Machine Learning Theory and Applications Laboratory \\ PUCRS
}
\date{\mespor de \the\year}

%----------------------------------------------------------------------------------------
%    EXTRA COLORS / COMMANDS FOR THIS LECTURE
%----------------------------------------------------------------------------------------
\definecolor{LightBlue}{HTML}{DBEAFE}
\definecolor{DarkBlue}{HTML}{1565C0}
\definecolor{DarkRed}{HTML}{C62828}
\definecolor{LightGray}{HTML}{F5F5F5}
\definecolor{MidGray}{HTML}{757575}
\definecolor{CellGreen}{HTML}{C8E6C9}
\definecolor{CellYellow}{HTML}{FFF9C4}
\definecolor{CellRed}{HTML}{FFCDD2}

%----------------------------------------------------------------------------------------
%    PRESENTATION SLIDES
%----------------------------------------------------------------------------------------

\begin{document}
  \begin{frame}
    \titlepage
  \end{frame}

  \begin{frame}{Visão Geral}
    \tableofcontents
  \end{frame}

  %======================================================================================
  % SECTION 0: INTRODUCTION & MOTIVATION
  %======================================================================================
  \section{Introdução e Motivação}

  %--- Slide: Why Position Matters ---
  \begin{frame}{Por que a posição importa?}
    \begin{columns}
      \begin{column}{0.52\textwidth}
        O mecanismo de \textit{self-attention} calcula \textit{scores} entre todos os pares:
        \begin{block}{Score de atenção}
          $\text{score}(i,j) = \mathbf{q}_i^\top \mathbf{k}_j \;/\; \sqrt{d}$
        \end{block}
        O produto escalar \textbf{não depende} dos índices $i$ ou $j$, apenas dos \textbf{valores} de $\mathbf{q}$ e $\mathbf{k}$.

        \vspace{0.3em}
        \textcolor{DarkRed}{\textbf{A \textit{self-attention} (bidirecional) é invariante a permutação.}}

        \vspace{0.3em}
        Sem informação posicional:\\
        ``\texttt{O gato sentou no tapete}'' $\equiv$ ``\texttt{O tapete sentou no gato}''
      \end{column}
      \begin{column}{0.45\textwidth}
        \textcolor{DarkBlue}{\textbf{O problema da permutação}}
        \vspace{0.3em}

        \resizebox{\linewidth}{!}{%
        \begin{tabular}{lll}
          \toprule
          \textbf{Token} & \textbf{Sem PE} & \textbf{Com PE} \\
          \midrule
          ``O''    & mesma repr. & pos 0 $\to$ única \\
          ``gato'' & mesma repr. & pos 1 $\to$ única \\
          ``sentou''& mesma repr. & pos 2 $\to$ única \\
          ``no''   & mesma repr. & pos 3 $\to$ única \\
          ``tapete''& mesma repr. & pos 4 $\to$ única \\
          \bottomrule
        \end{tabular}}
      \end{column}
    \end{columns}
  \end{frame}

  %--- Slide: Taxonomy ---
  \begin{frame}{Taxonomia de \textit{Positional Encodings}}
    \begin{columns}[t]
      \begin{column}{0.32\textwidth}
        \begin{block}{\textit{Absolute PE}}
          \begin{itemize}
            \item Vetor único por posição absoluta
            \item Sinusoidal (Vaswani et al., 2017)
            \item \textit{Learned embeddings} (BERT, GPT)
            \item Somado ao \textit{token embedding}
          \end{itemize}
        \end{block}
      \end{column}
      \begin{column}{0.32\textwidth}
        \begin{block}{\textit{Relative PE}}
          \begin{itemize}
            \item Codifica distância $|i - j|$
            \item RoPE: rotação no espaço
            \item ALiBi: \textit{bias} linear
            \item NoPE: máscara causal
          \end{itemize}
        \end{block}
      \end{column}
      \begin{column}{0.32\textwidth}
        \begin{block}{Extensão de Contexto}
          \begin{itemize}
            \item \textit{Fine-tune} para $L > L_{\text{train}}$
            \item \textit{Position Interpolation}
            \item YaRN
            \item LongRoPE
          \end{itemize}
        \end{block}
      \end{column}
    \end{columns}
  \end{frame}

  %--- Slide: Quadratic cost ---
  \begin{frame}{Comprimento de Contexto e o Custo Quadrático}
    \begin{columns}
      \begin{column}{0.5\textwidth}
        A atenção computa \textcolor{DarkBlue}{\textbf{todos os $L \times L$ pares}} $\to$ memória e computação escalam como $\mathbf{O(L^2)}$.

        \begin{itemize}
          \item Dobrar $L$ $\to$ 4$\times$ mais memória para a matriz de atenção
          \item \textcolor{DarkRed}{Um contexto de 100K \textit{tokens} em um modelo 7B requer dezenas de GB só de atenção (em \textit{fp16}, 2 bytes por valor)}
          \item Maioria dos modelos treina com $L = \text{2K--8K}$ \textit{tokens}
          \item Contextos longos desaceleram cada \textit{forward pass} quadraticamente
        \end{itemize}
      \end{column}
      \begin{column}{0.48\textwidth}
        \begin{tikzpicture}
          \begin{axis}[
            title={Memória vs.\ $L$ (escala $L^2$)},
            xlabel={Comprimento de sequência},
            ylabel={Memória relativa},
            symbolic x coords={1K,2K,4K,8K,16K,32K},
            xtick=data,
            ytick={0,200,400,600,800,1000},
            ymajorgrids=true,
            grid style=dashed,
            width=\linewidth,
            height=5cm,
            mark=*,
            thick,
            color=DarkBlue,
          ]
            \addplot coordinates {(1K,1)(2K,4)(4K,16)(8K,64)(16K,256)(32K,1024)};
          \end{axis}
        \end{tikzpicture}
        {\scriptsize\textit{Ilustrativo, memória de atenção escala como $O(L^2)$}}
      \end{column}
    \end{columns}
  \end{frame}

  %--- Slide: Economics ---
  \begin{frame}{Por que Extrapolação Importa: O Argumento Econômico}
    \textbf{Extrapolação:} usar o modelo em sequências \emph{mais longas} que as vistas no treino ($L_{\text{test}} > L_{\text{train}}$).

    \vspace{0.3em}
    \begin{itemize}
      \item \textbf{O custo de treinamento é dominado pela atenção quadrática}
        \begin{itemize}
          \item Treinar com $L = \text{2K}$ custa $4\times$ menos que $L = \text{4K}$
          \item Treinar com $L = \text{8K}$ custa $16\times$ menos que $L = \text{32K}$
        \end{itemize}
      \item \textbf{Se o PE puder extrapolar para sequências mais longas:}
        \begin{itemize}
          \item \textcolor{DarkBlue}{\textbf{Treinamento mais leve com $L$ pequeno, inferência com $L$ maior sem perda de desempenho}}
          \item Sem necessidade de \textit{fine-tuning} de contexto longo
        \end{itemize}
      \item \textcolor{DarkRed}{\textbf{Problema: a maioria dos PEs absolutos NÃO extrapola bem}}
        \begin{itemize}
          \item Esta aula: por que falham e como PEs relativos / extensão de contexto corrigem isso
        \end{itemize}
    \end{itemize}
  \end{frame}

  %--- Slide: Extrapolation challenge ---
  \begin{frame}{O Desafio da Extrapolação}
    Um modelo treinado com $L_{\text{train}} = \text{2K}$ encontra índices de posição \textbf{nunca vistos} quando $\textcolor{DarkRed}{L_{\text{test}} > L_{\text{train}}}$.

    \begin{columns}
      \begin{column}{0.55\textwidth}
        \vspace{0.3em}
        Perplexidade dispara além do horizonte de treino:
        \begin{tikzpicture}
          \begin{axis}[
            xlabel={Comprimento de contexto},
            ylabel={Perplexidade},
            symbolic x coords={512,1K,2K,3K,4K,6K,8K},
            xtick=data,
            ymajorgrids=true,
            grid style=dashed,
            width=\linewidth,
            height=4.5cm,
            mark=*,
            thick,
            color=DarkRed,
          ]
            \addplot coordinates {(512,12)(1K,12.5)(2K,13)(3K,20)(4K,55)(6K,140)(8K,310)};
          \end{axis}
        \end{tikzpicture}

        {\scriptsize\textit{Ilustrativo, comportamento típico além do horizonte de treino}}
      \end{column}
      \begin{column}{0.4\textwidth}
        \textbf{Como resolver:}
        \begin{itemize}
          \item PE que codifique apenas distância relativa $\to$ Parte 2
          \item Reescalar/interpolar posições em tempo de teste $\to$ Parte 4
          \item Sem PE: deixar o modelo inferir $\to$ NoPE
        \end{itemize}
      \end{column}
    \end{columns}
  \end{frame}

  %--- Slide: Notation ---
  \begin{frame}{Notação e Configuração}
    \centering
    \begin{tabular}{cl}
      \toprule
      \textbf{Símbolo} & \textbf{Significado} \\
      \midrule
      $d$ & Dimensão do \textit{embedding} ($d_{\text{model}}$) \\
      $L$ & Comprimento da sequência \\
      $i$ & Posição do \textit{token query} ($0 \le i < L$) \\
      $j$ & Posição do \textit{token key} ($0 \le j < L$) \\
      $k$ & Índice de dimensão ($0 \le k < d/2$) \\
      $\mathbf{x}_i$ & \textit{Embedding} de entrada do \textit{token} na posição $i$ \\
      $\mathbf{q}_i, \mathbf{k}_j$ & Vetores \textit{query} e \textit{key} nas posições $i$, $j$ \\
      $\text{PE}(i)$ & Vetor de \textit{positional encoding} para posição $i$ \\
      \bottomrule
    \end{tabular}

    \vspace{0.5em}
    \begin{block}{PE Absoluto}
      $\mathbf{x}_i \;\leftarrow\; \mathbf{x}_i + \text{PE}(i)$ \quad (somado ao \textit{embedding} antes da atenção)
    \end{block}
  \end{frame}

  %======================================================================================
  % SECTION 1: ABSOLUTE POSITIONAL ENCODING
  %======================================================================================
  \section{Absolute Positional Encoding}

  %--- Slide: Sinusoidal Paper ---
  \begin{frame}{PE Sinusoidal: Artigo}
    \begin{block}{Referência}
      Vaswani, A. et al. ``Attention Is All You Need.'' NeurIPS 2017.\footfullcite{vaswani2017attention}
    \end{block}
    \begin{itemize}
      \item Introduziu a arquitetura \textit{Transformer} (\textit{encoder--decoder})
      \item Propôs funções sinusoidais determinísticas como \textit{positional encodings}
      \item Propriedade: nenhum parâmetro aprendido $\to$ custo extra zero no treino
      \item Em princípio, pode generalizar para comprimentos não vistos no treino
      \item Tornou-se o PE padrão para a maioria das variantes iniciais do \textit{Transformer}
    \end{itemize}
  \end{frame}

  %--- Slide: Sinusoidal Motivation ---
  \begin{frame}{PE Sinusoidal: Por que Senoides?}
    \textbf{Queremos uma ``assinatura'' de posição que} (1) seja única para cada posição e (2) permita ao modelo recuperar \textit{distâncias relativas} entre \textit{tokens}.

    \vspace{0.4em}
    \begin{columns}[T]
      \begin{column}{0.52\textwidth}
        \begin{block}{Ideia: sobrepor ondas de várias frequências}
          Como na contagem binária, cada \textit{bit} alterna a uma taxa diferente:
          \begin{itemize}
            \item Componentes \textbf{rápidos} (alta freq.) distinguem posições \textbf{vizinhas}
            \item Componentes \textbf{lentos} (baixa freq.) capturam a posição \textbf{global}
          \end{itemize}
          Senoides são o análogo \textbf{contínuo} dos \textit{bits}.
        \end{block}
      \end{column}
      \begin{column}{0.44\textwidth}
        \begin{exampleblock}{Termos}
          \textbf{Frequência} $\omega$: quão rápido a onda oscila.\\[2pt]
          \textbf{Período} (comprimento de onda) $=2\pi/\omega$: a distância para a onda se repetir.
        \end{exampleblock}
        \vspace{0.2em}
        {\small Por que seno \emph{e} cosseno? Um deslocamento de $\delta$ vira uma \textbf{rotação} do par $(\sin,\cos)$, é isso que deixa a atenção ler \textit{offsets} relativos (próximos slides).}
      \end{column}
    \end{columns}
  \end{frame}

  %--- Slide: Sinusoidal Formula ---
  \begin{frame}{PE Sinusoidal: Fórmula}
    Para posição $i$ e índice de dimensão $k$ ($d$ = dimensão total do \textit{embedding}):
    \begin{block}{Fórmula do \textit{Positional Encoding} Sinusoidal}
      \begin{align*}
        \text{PE}(i, 2k)   &= \sin\!\left( \frac{i}{10000^{2k/d}} \right) \\[4pt]
        \text{PE}(i, 2k{+}1) &= \cos\!\left( \frac{i}{10000^{2k/d}} \right)
      \end{align*}
    \end{block}
    \begin{itemize}
      \item Dimensão $2k$ usa \textbf{seno}; dimensão $2k{+}1$ usa \textbf{cosseno}, mesma frequência, defasagem de 90°
      \item Base $10000$ (empírica): frequências altas (dim 0) a baixas (dim $d{-}1$)
    \end{itemize}
    \begin{exampleblock}{Abreviação}
      $\omega_k = 10000^{-2k/d}$ \quad$\Rightarrow$\quad $\text{PE}(i, 2k) = \sin(i \cdot \omega_k)$, \; $\text{PE}(i, 2k{+}1) = \cos(i \cdot \omega_k)$
    \end{exampleblock}
  \end{frame}

  %--- Slide: Sinusoidal Intuition ---
  \begin{frame}{PE Sinusoidal: Intuição}
    \begin{columns}
      \begin{column}{0.5\textwidth}
        \textcolor{DarkBlue}{\textbf{Analogia com contagem binária:}}

        \vspace{0.3em}
        \resizebox{\linewidth}{!}{%
        \begin{tabular}{cccc}
          \toprule
          \textbf{pos} & bit 2 (lento) & bit 1 (médio) & bit 0 (rápido) \\
          \midrule
          0 & 0 & 0 & 0 \\
          1 & 0 & 0 & 1 \\
          2 & 0 & 1 & 0 \\
          3 & 0 & 1 & 1 \\
          4 & 1 & 0 & 0 \\
          5 & 1 & 0 & 1 \\
          \bottomrule
        \end{tabular}}

        \vspace{0.5em}
        $\to$ PE sinusoidal é o análogo contínuo.
      \end{column}
      \begin{column}{0.47\textwidth}
        {\small
        Cada dimensão = um ``\textit{canal de frequência}''

        \vspace{0.3em}
        Low-$k$: oscilação rápida $\to$ posição fina\\
        High-$k$: oscilação lenta $\to$ posição grosseira

        \vspace{0.4em}
        \textcolor{DarkRed}{\textbf{Nomenclatura:} \textbf{low-$k$} = índice $k$ \textit{pequeno} = \textbf{alta} frequência; \textbf{high-$k$} = índice $k$ \textit{grande} = \textbf{baixa} frequência.}
        }

        \vspace{0.5em}
        $\text{PE}(i+\delta)$ pode ser expresso como função linear de $\text{PE}(i)$, permitindo ao modelo aprender \textit{offsets} relativos a partir de PEs absolutos \textit{(é a propriedade-chave; veja o slide de Propriedades)}.
      \end{column}
    \end{columns}
  \end{frame}

  %--- Slide: Sinusoidal Multi-Frequency Waves Visualization ---
  \begin{frame}{PE Sinusoidal: Visualizando as Frequências ($d=16$)}
    \centering
    \begin{tikzpicture}
      \begin{axis}[
        width=11cm, height=5.8cm,
        xlabel={posição $i$},
        ylabel={$\sin(i \cdot \omega_k)$},
        xmin=0, xmax=30, ymin=-1.2, ymax=1.2,
        grid=major, grid style={dashed, gray!30},
        ymajorgrids=true,
        legend style={font=\scriptsize, at={(0.98,1.02)}, anchor=south east, legend columns=4, draw=none, fill=none, /tikz/every even column/.append style={column sep=0.3cm}},
        label style={font=\footnotesize}, tick label style={font=\scriptsize},
      ]
        % k=0: fastest (omega = 1.0)
        \addplot[DarkBlue, very thick, domain=0:30, samples=200] {sin(deg(x*1.0))};
        \addlegendentry{$k{=}0$ (low-$k$, rápido)}

        % k=1: medium-fast
        \addplot[DarkRed, very thick, domain=0:30, samples=200] {sin(deg(x*0.316))};
        \addlegendentry{$k{=}1$}

        % k=2: medium-slow
        \addplot[MidGray, very thick, domain=0:30, samples=200] {sin(deg(x*0.1))};
        \addlegendentry{$k{=}2$}

        % k=3: slowest (omega ~ 0.032)
        \addplot[ForestGreen, very thick, domain=0:30, samples=200] {sin(deg(x*0.0316))};
        \addlegendentry{$k{=}3$ (high-$k$, lento)}
      \end{axis}
    \end{tikzpicture}

    \vspace{0.1em}
    {\scriptsize
    Cada $k$ é um ``relógio'' em velocidade diferente: combinando frequências, cada posição tem uma assinatura única (análogo contínuo da contagem binária). \textcolor{DarkRed}{\textbf{A seguir:}} frase ``\textit{O gato sentou no tapete}'' ($L{=}5$, $d{=}4$).
    }
  \end{frame}

  %--- Slide: Sinusoidal Visualisation (Preview of Hand-Calc Target) ---
  \begin{frame}{PE Sinusoidal: Matriz para ``O gato sentou no tapete''}
    Matriz PE completa para a frase exemplo ($L=5$, $d=4$):

    \vspace{0.3em}
    \centering
    \begin{tabular}{c|rrrr}
      \toprule
      \textbf{$i$ (token)} & \textbf{dim 0} & \textbf{dim 1} & \textbf{dim 2} & \textbf{dim 3} \\
      \midrule
      \rowcolor{CellGreen} 0 (``O'')      &  0.000 &  1.000 & 0.000 & 1.000 \\
      \rowcolor{CellGreen} 1 (``gato'')   &  0.841 &  0.540 & 0.010 & 1.000 \\
      \rowcolor{CellGreen} 2 (``sentou'') &  0.909 & $-$0.416 & 0.020 & 1.000 \\
      \rowcolor{CellGreen} 3 (``no'')     &  0.141 & $-$0.990 & 0.030 & 1.000 \\
      \rowcolor{CellGreen} 4 (``tapete'') & $-$0.757 & $-$0.654 & 0.040 & 0.999 \\
      \bottomrule
    \end{tabular}

    \vspace{0.5em}
    \flushleft
    Dimensões 0--1 ($\omega_0{=}1$) variam rapidamente com $i$; dimensões 2--3 ($\omega_1{=}0.01$) quase não mudam nessa faixa.\\
    \textcolor{DarkBlue}{\textbf{Os próximos slides mostram passo a passo de onde vêm esses números.}}
  \end{frame}

  %--- Slide: Hand Computation Step 1 ---
  \begin{frame}{Cálculo Manual, PE Sinusoidal, Passo 1: Configuração}
    \textbf{Frase:} ``O gato sentou no tapete'' $\to$ $L = 5$, posições $i \in \{0, 1, 2, 3, 4\}$, $d = 4$

    \vspace{0.3em}
    Calcular as frequências angulares $\omega_k = 10000^{-2k/d}$:

    \vspace{0.3em}
    \centering
    \begin{tabular}{ccccc}
      \toprule
      $k$ & $2k/d$ & Expoente & $\omega_k$ & Período $2\pi/\omega_k$ \\
      \midrule
      \rowcolor{CellGreen} 0 & $0/4 = 0.0$  & $10000^0 = 1$   & $\omega_0 = 1.0000$ & $\approx 6.28$ \\
      \rowcolor{CellGreen} 1 & $2/4 = 0.5$  & $10000^{0.5} = 100$ & $\omega_1 = 0.0100$ & $\approx 628$ \\
      \bottomrule
    \end{tabular}

    \vspace{0.5em}
    \flushleft
    Com $d=4$ temos 2 pares de frequência ($k = 0, 1$).\\
    $k{=}0$: oscilação rápida (período $\approx$ 6 tokens) $\;|\;$ $k{=}1$: muito lenta (período $\approx$ 628 tokens).
  \end{frame}

  %--- Slide: Hand Computation Step 2 ---
  \begin{frame}{Cálculo Manual, PE Sinusoidal, Passo 2: Preencher a Tabela}
    Aplicar $\text{PE}(i, 2k) = \sin(i \cdot \omega_k)$ e $\text{PE}(i, 2k{+}1) = \cos(i \cdot \omega_k)$ aos 5 tokens:

    \vspace{0.3em}
    \resizebox{\linewidth}{!}{%
    \begin{tabular}{ccccc}
      \toprule
      $i$ (token) & dim 0 \; $\sin(i \cdot \omega_0)$ & dim 1 \; $\cos(i \cdot \omega_0)$ & dim 2 \; $\sin(i \cdot \omega_1)$ & dim 3 \; $\cos(i \cdot \omega_1)$ \\
      \midrule
      \rowcolor{CellGreen} 0 (``O'')      & $\sin(0) = 0.000$  & $\cos(0) = 1.000$      & $\sin(0) = 0.000$     & $\cos(0) = 1.000$ \\
      \rowcolor{CellGreen} 1 (``gato'')   & $\sin(1) = 0.841$  & $\cos(1) = 0.540$      & $\sin(0.01) = 0.010$  & $\cos(0.01) = 1.000$ \\
      \rowcolor{CellGreen} 2 (``sentou'') & $\sin(2) = 0.909$  & $\cos(2) = {-}0.416$   & $\sin(0.02) = 0.020$  & $\cos(0.02) = 1.000$ \\
      \rowcolor{CellGreen} 3 (``no'')     & $\sin(3) = 0.141$  & $\cos(3) = {-}0.990$   & $\sin(0.03) = 0.030$  & $\cos(0.03) = 1.000$ \\
      \rowcolor{CellGreen} 4 (``tapete'') & $\sin(4) = {-}0.757$ & $\cos(4) = {-}0.654$ & $\sin(0.04) = 0.040$  & $\cos(0.04) = 0.999$ \\
      \bottomrule
    \end{tabular}}

    \vspace{0.5em}
    Os valores reproduzem exatamente a matriz prevista no slide anterior. As dimensões 2 e 3 quase não mudam: o modelo usa essas dimensões para codificar posição \textbf{grosseira} (milhares de passos).
  \end{frame}

  %--- Slide: Hand Computation Step 3 ---
  \begin{frame}{Cálculo Manual, Passo 3: Somar PE ao \textit{Embedding}}
    {\footnotesize \textit{Token embeddings} de ``O gato sentou no tapete'':}

    \vspace{0.1em}
    {\scriptsize
    \centering
    \begin{tabular}{ccccc}
      \toprule
      $i$ (token) & $x[0]$ & $x[1]$ & $x[2]$ & $x[3]$ \\
      \midrule
      0 (``O'')      &  0.500 & $-$0.200 &  0.100 &  0.800 \\
      1 (``gato'')   &  0.300 &  0.600 & $-$0.400 &  0.200 \\
      2 (``sentou'') & $-$0.100 &  0.400 &  0.700 & $-$0.300 \\
      3 (``no'')     &  0.600 & $-$0.500 &  0.200 &  0.100 \\
      4 (``tapete'') & $-$0.200 &  0.300 & $-$0.600 &  0.500 \\
      \bottomrule
    \end{tabular}}

    \vspace{0.2em}
    {\footnotesize Após somar o PE sinusoidal:}

    \vspace{0.1em}
    \resizebox{\linewidth}{!}{%
    \begin{tabular}{ccccc}
      \toprule
      $i$ & $x[0]+\text{PE}[0]$ & $x[1]+\text{PE}[1]$ & $x[2]+\text{PE}[2]$ & $x[3]+\text{PE}[3]$ \\
      \midrule
      \rowcolor{CellGreen} 0 & $0.500+0.000=0.500$   & $-0.200+1.000=0.800$   & $0.100+0.000=0.100$   & $0.800+1.000=1.800$ \\
      \rowcolor{CellGreen} 1 & $0.300+0.841=1.141$   & $0.600+0.540=1.140$    & $-0.400+0.010=-0.390$ & $0.200+1.000=1.200$ \\
      \rowcolor{CellGreen} 2 & $-0.100+0.909=0.809$  & $0.400-0.416=-0.016$   & $0.700+0.020=0.720$   & $-0.300+1.000=0.700$ \\
      \rowcolor{CellGreen} 3 & $0.600+0.141=0.741$   & $-0.500-0.990=-1.490$  & $0.200+0.030=0.230$   & $0.100+1.000=1.100$ \\
      \rowcolor{CellGreen} 4 & $-0.200-0.757=-0.957$ & $0.300-0.654=-0.354$   & $-0.600+0.040=-0.560$ & $0.500+0.999=1.499$ \\
      \bottomrule
    \end{tabular}}

    \vspace{0.2em}
    {\scriptsize \textcolor{DarkBlue}{$\checkmark$ A mesma rede processa os 5 \textit{tokens}; a posição fica codificada nos valores.}}
  \end{frame}

  %--- Slide: Properties ---
  \begin{frame}{PE Sinusoidal: Propriedades}
    \begin{itemize}
      \item \textbf{Determinístico}: nenhum parâmetro aprendido, idêntico entre execuções
      \item \textbf{Limitado}: todos os valores em $[-1, 1]$, estável com normalização
      \item \textbf{Único por posição}: cada $i$ gera um vetor PE distinto
      \item \textcolor{DarkBlue}{\textbf{\textit{Offsets} relativos são lineares}}: $\text{PE}(i+\delta) = M(\delta) \cdot \text{PE}(i)$ para uma matriz fixa $M$
        \begin{itemize}
          \item Permite ao modelo aprender distância relativa a partir de PE absoluto
        \end{itemize}
      \item \textcolor{DarkBlue}{\textbf{Teoricamente generalizável}}: posições além do treino existem na mesma curva
        \begin{itemize}
          \item \textcolor{DarkRed}{Na prática, a generalização é limitada $\to$ ver próximo slide}
        \end{itemize}
    \end{itemize}
  \end{frame}

  %--- Slide: Learned Absolute PE ---
  \begin{frame}{\textit{Learned Absolute PE}}
    \begin{itemize}
      \item \textbf{BERT, GPT iniciais:} PE é uma matriz de \textit{embeddings} treinável $E \in \mathbb{R}^{L_{\max} \times d}$
      \item Consulta por posição: $\text{PE}(i) = E[i, :]$
      \item \textbf{Vantagens:} adapta-se à distribuição dos dados; sem frequências manuais
      \item \textbf{Desvantagens:}
        \begin{itemize}
          \item \textcolor{DarkRed}{Limite rígido em $L_{\max}$: posições além simplesmente não existem}
          \item \textcolor{DarkRed}{\textit{Overfits} ao comprimento de treino; pior OOD (\textit{out-of-distribution}, fora da distribuição de treino) que sinusoidal}
          \item Adiciona $L_{\max} \times d$ parâmetros ao modelo
        \end{itemize}
      \item \textcolor{DarkBlue}{\textbf{Modelos modernos migraram para PE relativo (RoPE, ALiBi) $\to$}}
    \end{itemize}
  \end{frame}

  %--- Slide: Limitation ---
  \begin{frame}{Limitação do PE Absoluto}
    \begin{itemize}
      \item O modelo vê posições absolutas durante o treino: $i = 0, 1, \ldots, L_{\text{train}}{-}1$
      \item \textcolor{DarkRed}{\textbf{Posições $i \ge L_{\text{train}}$ são \textit{out-of-distribution}}}: nunca vistas no treino
      \item As frequências sinusoidais não extrapolam os padrões de atenção de forma confiável
      \item \textcolor{DarkRed}{\textit{Learned absolute} PE (BERT) é ainda pior, sem forma fechada para novas posições}
      \item \textbf{Problema central: posição codificada independentemente do contexto}
        \begin{itemize}
          \item ``gato'' na posição 0 e na posição 100 recebem vetores PE diferentes, mesmo que a relação semântica seja a mesma
        \end{itemize}
      \item \textcolor{DarkBlue}{\textbf{Isso motiva os \textit{positional encodings} relativos $\to$}}
    \end{itemize}
  \end{frame}

  %--- Slide: Transition ---
  \begin{frame}{De Absoluto para Relativo}
    \textcolor{DarkRed}{\textbf{O problema fundamental do PE absoluto:}}

    \begin{alertblock}{}
      Dois \textit{tokens} ``gato'' à distância 2 de ``sentou'' deveriam ter o mesmo \textit{score} de atenção,
      independentemente de estarem nas posições $(0, 2)$ ou $(1000, 1002)$.
    \end{alertblock}

    \textcolor{DarkBlue}{\textbf{PE relativo codifica $j - i$ diretamente:}}

    \begin{itemize}
      \item O \textit{score} de atenção depende apenas do \textit{offset} relativo $j - i$, não de $i$ ou $j$ individualmente
      \item Modelos treinados com comprimento $L$ podem generalizar para qualquer comprimento
      \item Três estratégias: \textbf{RoPE} (rotação), \textbf{ALiBi} (\textit{bias} aditivo), \textbf{NoPE} (implícito)
    \end{itemize}
  \end{frame}

  %======================================================================================
  % SECTION 2: RELATIVE POSITIONAL ENCODING
  %======================================================================================
  \section{Relative Positional Encoding}

  %--- Slide: Key Idea ---
  \begin{frame}{Ideia Central: \textit{Positional Encodings} Relativos}
    \textcolor{DarkBlue}{\textbf{Objetivo: fazer o \textit{score} de atenção depender apenas de $(j - i)$.}}

    \begin{block}{Formulação}
      $\text{score}(i, j) = f(\mathbf{q}_i, \mathbf{k}_j, j{-}i)$ \quad onde $f$ depende de $(j{-}i)$, não de $i$ ou $j$ separadamente
    \end{block}

    \vspace{0.3em}
    Três estratégias:
    \begin{itemize}
      \item \textcolor{DarkBlue}{\textbf{RoPE}}: rotacionar $\mathbf{q}$ e $\mathbf{k}$ por ângulo dependente da posição; produto escalar codifica $j{-}i$
      \item \textcolor{DarkBlue}{\textbf{ALiBi}}: somar penalidade linear $-m \cdot |i{-}j|$ aos \textit{logits} de atenção
      \item \textcolor{DarkBlue}{\textbf{NoPE}}: sem PE; posição implícita via máscara causal
    \end{itemize}

    Todas alcançam generalização de comprimento em graus variados.
  \end{frame}

  %--- SUBSECTION: RoPE ---
  \subsection{RoPE}

  %--- Slide: RoPE Paper ---
  \begin{frame}{RoPE: Artigo}
    \begin{block}{Referência}
      Su, J. et al. ``RoFormer: Enhanced Transformer with Rotary Position Embedding.'' Neurocomputing 2024.\footfullcite{su2024roformer}
    \end{block}
    \begin{itemize}
      \item Codificar posição \textbf{rotacionando} os vetores \textit{query} e \textit{key}
      \item Rotação é operação linear, compatível com qualquer arquitetura \textit{Transformer}
      \item O produto interno $\mathbf{q}_i^\top \mathbf{k}_j$ depende apenas de $(j - i)$ após a rotação $\checkmark$
      \item Adotado por LLaMA, Mistral, Falcon, Gemma e maioria dos LLMs \textit{open-source}
      \item Base para métodos de extensão de contexto (YaRN, LongRoPE, etc.)
    \end{itemize}
  \end{frame}

  %--- Slide: RoPE 2D Rotation Diagram ---
  \begin{frame}{RoPE: Visualizando a Rotação 2D}
    \textbf{Intuição do RoPE: posição $=$ ângulo de rotação} (a álgebra vem a seguir).

    \vspace{0.2em}
    \begin{columns}
      \begin{column}{0.45\textwidth}
        {\small Bloco $k{=}0$: rotacionar $\mathbf{q} = (1, 0)$ por $i \cdot \theta_0 = 1.0$ rad}

        \centering
        \begin{tikzpicture}[scale=1.5]
          % Axes
          \draw[->, gray] (-0.3,0) -- (1.3,0) node[right, font=\tiny] {$q_0$};
          \draw[->, gray] (0,-0.3) -- (0,1.3) node[above, font=\tiny] {$q_1$};
          % Unit circle
          \draw[gray, dashed, thin] (0,0) circle (1);
          % Original vector (blue)
          \draw[->, DarkBlue, very thick] (0,0) -- (1,0) node[below right, font=\small] {\textcolor{DarkBlue}{$\mathbf{q}$}};
          % Rotated vector (red)
          \draw[->, DarkRed, very thick] (0,0) -- (0.540,0.841) node[above right, font=\small] {\textcolor{DarkRed}{$\tilde{\mathbf{q}}$}};
          % Angle arc
          \draw[thick] (0.35,0) arc (0:57.3:0.35);
          \node[font=\scriptsize] at (0.55,0.3) {$57.3^{\circ}$};
          % Projections (dashed)
          \draw[gray, dashed, thin] (0.540,0.841) -- (0.540,0);
          \draw[gray, dashed, thin] (0.540,0.841) -- (0,0.841);
          % Projection labels
          \node[below, font=\tiny, gray] at (0.540,0) {$0.540$};
          \node[left, font=\tiny, gray] at (0,0.841) {$0.841$};
        \end{tikzpicture}
      \end{column}
      \begin{column}{0.52\textwidth}
        \textbf{Bloco $k{=}0$: \textit{query} em $i=1$}

        \begin{exampleblock}{Multiplicação matricial}
          $\begin{pmatrix} \cos 1.0 & -\sin 1.0 \\ \sin 1.0 & \cos 1.0 \end{pmatrix} \begin{pmatrix} 1 \\ 0 \end{pmatrix} = \begin{pmatrix} 0.540 \\ 0.841 \end{pmatrix}$
        \end{exampleblock}

        \begin{itemize}
          \item O vetor é rotacionado 57.3° no sentido anti-horário
          \item \textcolor{DarkBlue}{\textbf{Rotação preserva a magnitude: $|\mathbf{q}| = |\tilde{\mathbf{q}}| = 1$ $\checkmark$}}
          \item \textcolor{DarkBlue}{$\tilde{\mathbf{q}}_i^\top \tilde{\mathbf{k}}_j$ passará a depender só de $(j{-}i) \cdot \theta_0$ \textit{(provado a seguir)}}
        \end{itemize}
      \end{column}
    \end{columns}
  \end{frame}

  %--- Slide: RoPE Core Idea ---
  \begin{frame}{RoPE: Ideia Central}
    Em vez de somar PE ao \textit{embedding}, \textbf{rotacionar Q e K}:

    \begin{block}{Rotação e \textit{Score}}
      \vspace{-0.5em}
      \begin{align*}
        \tilde{\mathbf{q}}_i &= R_\theta(i) \cdot \mathbf{q}_i \qquad\qquad \tilde{\mathbf{k}}_j = R_\theta(j) \cdot \mathbf{k}_j \\[4pt]
        \text{score}(i,j) &= \tilde{\mathbf{q}}_i^\top \tilde{\mathbf{k}}_j = \mathbf{q}_i^\top R_\theta(i)^\top R_\theta(j)\, \mathbf{k}_j = \mathbf{q}_i^\top R_\theta(j{-}i)\, \mathbf{k}_j
      \end{align*}
    \end{block}

    Como matrizes de rotação satisfazem $R_\theta(i)^\top = R_\theta(-i)$ (rotação no sentido oposto), temos $R_\theta(i)^\top R_\theta(j) = R_\theta(-i)\,R_\theta(j) = R_\theta(j{-}i)$, \textcolor{DarkBlue}{\textbf{o \textit{score} depende apenas do \textit{offset} relativo $j - i$.}}

    \begin{itemize}
      \item $R_\theta(m)$ é uma matriz de rotação bloco-diagonal parametrizada pela posição $m$
      \item Cada bloco 2D rotaciona um par de dimensões por ângulo $m \cdot \theta_k$
      \item \textcolor{DarkBlue}{\textbf{PE NÃO é somado à entrada, é aplicado dentro do \textit{attention head}}}
    \end{itemize}
  \end{frame}

  %--- Slide: RoPE 2D Block ---
  \begin{frame}{RoPE: O Bloco de Rotação 2D}
    Para um par de dimensões $(q_0, q_1)$ na posição $i$:

    \begin{block}{Matriz de rotação 2D}
      $R_\theta(i) = \begin{pmatrix} \cos(i\theta) & -\sin(i\theta) \\ \sin(i\theta) & \phantom{-}\cos(i\theta) \end{pmatrix}$

      \vspace{0.3em}
      Par rotacionado: $\bigl[\cos(i\theta)\,q_0 - \sin(i\theta)\,q_1,\;\; \sin(i\theta)\,q_0 + \cos(i\theta)\,q_1\bigr]$
    \end{block}

    A mesma rotação é aplicada ao par correspondente de \textit{key}, no ângulo $j \cdot \theta$.

    \vspace{0.3em}
    \textcolor{DarkBlue}{\textbf{O produto escalar equivale a uma rotação por $(j{-}i) \cdot \theta$, depende só do \textit{offset}.}}

    \begin{exampleblock}{Expansão}
      Somando os dois produtos componente a componente, os termos cruzados de $\sin$/$\cos$ recombinam-se em $\cos((j{-}i)\theta)$ e $\sin((j{-}i)\theta)$ $\to$ depende apenas de $(j-i)$. \textit{Derivação completa no próximo slide.}
    \end{exampleblock}
  \end{frame}

  %--- Slide: RoPE Full Formula ---
  \begin{frame}{RoPE: Fórmula Completa ($d$ Dimensões)}
    Para $d$ dimensões, aplicar $d/2$ blocos de rotação independentes:

    \begin{block}{Frequências e rotação bloco-diagonal}
      \vspace{-0.5em}
      \begin{align*}
        \theta_k &= 10000^{-2k/d} \qquad k = 0, 1, \ldots, d/2{-}1 \\[4pt]
        R_\theta(i) &= \text{block-diag}\!\bigl(R(i\theta_0),\; R(i\theta_1),\; \ldots,\; R(i\theta_{d/2-1})\bigr)
      \end{align*}
    \end{block}

    \begin{itemize}
      \item Mesma base $10000$ do PE sinusoidal: frequências cobrem muitas ordens de magnitude
      \item \textit{Low-k} ($\theta_k \approx 1$): rotação rápida $\to$ posição relativa fina
      \item \textit{High-k} ($\theta_k \approx 0$): rotação lenta $\to$ posição relativa grosseira
      \item Eficiente: tabelas cos/sin pré-computadas; fundido com projeção Q$\cdot$K
    \end{itemize}
  \end{frame}

  %--- Slide: RoPE Frequency Spectrum Visualization ---
  \begin{frame}{RoPE: Espectro de Frequências}
    \centering
    \begin{tikzpicture}
      \begin{semilogyaxis}[
        width=11cm, height=5.8cm,
        xlabel={índice da dimensão $k$ (com $d=64$)},
        ylabel={$\theta_k = 10000^{-2k/d}$},
        xmin=0, xmax=32, ymin=0.00005, ymax=2,
        grid=major, grid style={dashed, gray!30},
        ymajorgrids=true,
        label style={font=\footnotesize}, tick label style={font=\scriptsize},
        log basis y=10,
      ]
        % Shaded regions for low-k, mid-k, high-k
        \fill[DarkBlue, opacity=0.10] (axis cs:0,0.1) rectangle (axis cs:32,2);
        \fill[MidGray, opacity=0.15] (axis cs:0,0.001) rectangle (axis cs:32,0.1);
        \fill[DarkRed,  opacity=0.10] (axis cs:0,0.00005) rectangle (axis cs:32,0.001);

        % Labels for regions
        \node[font=\scriptsize, DarkBlue, anchor=east] at (axis cs:31,0.5)    {\textbf{low-$k$}: rotação rápida};
        \node[font=\scriptsize, MidGray, anchor=east] at (axis cs:31,0.01)   {\textbf{mid-$k$}};
        \node[font=\scriptsize, DarkRed,  anchor=east] at (axis cs:31,0.00015) {\textbf{high-$k$}: rotação lenta};

        % The curve
        \addplot[DarkBlue, very thick, mark=*, mark size=1.4pt, samples=33, domain=0:32]
          {10000^(-2*x/64)};
      \end{semilogyaxis}
    \end{tikzpicture}

    \vspace{0.2em}
    {\footnotesize
    Cobertura de $\sim$4 ordens de magnitude em frequência. \textcolor{DarkBlue}{\textbf{Diversidade espectral é a base do sucesso do RoPE}}: cada bloco 2D rotaciona numa velocidade diferente, permitindo codificar posições relativas finas e grosseiras simultaneamente.
    }
  \end{frame}

  %--- Slide: RoPE Proof Sketch ---
  \begin{frame}{RoPE, Por que Funciona: Esboço de Prova}
    \textbf{Afirmação:} $\tilde{\mathbf{q}}_i^\top \tilde{\mathbf{k}}_j$ depende apenas de $(j - i)$.

    \vspace{0.3em}
    Foco em um bloco 2D com frequência $\theta$:

    \begin{block}{Derivação}
      \vspace{-0.5em}
      \begin{align*}
        \tilde{\mathbf{q}}_i &= \bigl(\cos(i\theta)\,q_0 - \sin(i\theta)\,q_1,\;\; \sin(i\theta)\,q_0 + \cos(i\theta)\,q_1\bigr)\\[2pt]
        \tilde{\mathbf{k}}_j &= \bigl(\cos(j\theta)\,k_0 - \sin(j\theta)\,k_1,\;\; \sin(j\theta)\,k_0 + \cos(j\theta)\,k_1\bigr)\\[6pt]
        \tilde{\mathbf{q}}_i \cdot \tilde{\mathbf{k}}_j &= (q_0 k_0 + q_1 k_1)\cos\!\bigl((j{-}i)\theta\bigr) + (q_1 k_0 - q_0 k_1)\sin\!\bigl((j{-}i)\theta\bigr) \\[2pt]
        &\to \text{depende apenas de } j - i \;\; \checkmark
      \end{align*}
    \end{block}

    \begin{itemize}
      \item Cada bloco contribui um par $(\cos, \sin)$ que depende de $(j-i)$
      \item \textcolor{DarkBlue}{\textbf{Somando $d/2$ blocos: \textit{score} total é função de $(\mathbf{q}, \mathbf{k}, j{-}i)$}}
    \end{itemize}
  \end{frame}

  %--- Slide: RoPE Hand Computation Step 1 ---
  \begin{frame}{Cálculo Manual, RoPE, Passo 1: Configuração}
    \textbf{Parâmetros:} $d = 4$, \; $\theta_0 = 1.0$, \; $\theta_1 = 0.01$\\
    \textit{Query} em $i=1$: $\mathbf{q} = [1, 0, 1, 0]$ \qquad \textit{Key} em $j=3$: $\mathbf{k} = [0, 1, 0, 1]$

    \vspace{0.3em}
    Ângulos de rotação:

    {\small
    \begin{tabular}{cccc}
      \toprule
      Bloco & $\theta_k$ & Ângulo $i{=}1$ (\textit{query}) & Ângulo $j{=}3$ (\textit{key}) \\
      \midrule
      \rowcolor{CellGreen} $k{=}0$ & 1.0000 & $1 \times 1.00 = 1.0000$ rad & $3 \times 1.00 = 3.0000$ rad \\
      \rowcolor{CellGreen} $k{=}1$ & 0.0100 & $1 \times 0.01 = 0.0100$ rad & $3 \times 0.01 = 0.0300$ rad \\
      \bottomrule
    \end{tabular}}

    \vspace{0.3em}
    Valores trigonométricos:

    {\small
    \begin{tabular}{ccc}
      \toprule
      ângulo & $\cos$ & $\sin$ \\
      \midrule
      1.0000 & 0.540 & 0.841 \\
      3.0000 & $-$0.990 & 0.141 \\
      0.0100 & 1.000 & 0.010 \\
      0.0300 & 1.000 & 0.030 \\
      \bottomrule
    \end{tabular}}
  \end{frame}

  %--- Slide: RoPE Hand Computation Step 2 ---
  \begin{frame}{Cálculo Manual, RoPE, Passo 2: Aplicar Rotações}
    \textbf{Rotacionar $\mathbf{q} = [1, 0, 1, 0]$ em $i = 1$:}

    {\small
    \begin{tabular}{ccccc}
      \toprule
      Bloco & Entrada & Ângulo & Fórmula & Resultado \\
      \midrule
      \rowcolor{CellGreen} $k{=}0$ & $(1,0)$ & $1.00$ rad & $\cos \cdot 1 - \sin \cdot 0,\; \sin \cdot 1 + \cos \cdot 0$ & $(0.540, 0.841)$ \\
      \rowcolor{CellGreen} $k{=}1$ & $(1,0)$ & $0.01$ rad & $\cos \cdot 1 - \sin \cdot 0,\; \sin \cdot 1 + \cos \cdot 0$ & $(1.000, 0.010)$ \\
      \bottomrule
    \end{tabular}}

    $\to$ \textcolor{DarkBlue}{$\tilde{\mathbf{q}}(i{=}1) = [0.540,\; 0.841,\; 1.000,\; 0.010]$}

    \vspace{0.5em}
    \textbf{Rotacionar $\mathbf{k} = [0, 1, 0, 1]$ em $j = 3$:}

    {\small
    \begin{tabular}{ccccc}
      \toprule
      Bloco & Entrada & Ângulo & Fórmula & Resultado \\
      \midrule
      \rowcolor{CellGreen} $k{=}0$ & $(0,1)$ & $3.00$ rad & $\cos \cdot 0 - \sin \cdot 1,\; \sin \cdot 0 + \cos \cdot 1$ & $(-0.141, -0.990)$ \\
      \rowcolor{CellGreen} $k{=}1$ & $(0,1)$ & $0.03$ rad & $\cos \cdot 0 - \sin \cdot 1,\; \sin \cdot 0 + \cos \cdot 1$ & $(-0.030, 1.000)$ \\
      \bottomrule
    \end{tabular}}

    $\to$ \textcolor{DarkBlue}{$\tilde{\mathbf{k}}(j{=}3) = [-0.141,\; -0.990,\; -0.030,\; 1.000]$}
  \end{frame}

  %--- Slide: RoPE Step 3 ---
  \begin{frame}{Cálculo Manual, RoPE, Passo 3: Produto Escalar e Verificação}
    $\text{score}(i{=}1, j{=}3) = \tilde{\mathbf{q}}(1) \cdot \tilde{\mathbf{k}}(3)$:

    \begin{block}{Cálculo}
      $= 0.540 \!\cdot\! ({-}0.141) + 0.841 \!\cdot\! ({-}0.990) + 1.000 \!\cdot\! ({-}0.030) + 0.010 \!\cdot\! 1.000 = \mathbf{-0.929}$
    \end{block}

    \textcolor{DarkBlue}{\textbf{Verificação: depende apenas de $(j - i) = 2$?}}

    Rotacionar o mesmo $\mathbf{q}$ em $i=0$ e o mesmo $\mathbf{k}$ em $j=2$ (\textit{offset} relativo = 2):

    \begin{exampleblock}{}
      {\small
      $\tilde{\mathbf{q}}(0) = [1, 0, 1, 0]$ \; (rotação por 0: identidade)\\[2pt]
      $\tilde{\mathbf{k}}(2)$: bloco 0 $\to$ $(-0.909, -0.416)$; \; bloco 1 $\to$ $(-0.020, 1.000)$\\[2pt]
      $\tilde{\mathbf{k}}(2) = [-0.909, -0.416, -0.020, 1.000]$}
    \end{exampleblock}

    \begin{block}{Verificação}
      $\text{score} = 1 \!\cdot\! ({-}0.909) + 0 \!\cdot\! ({-}0.416) + 1 \!\cdot\! ({-}0.020) + 0 \!\cdot\! 1.000 = \mathbf{-0.929}$ \;\; $\checkmark$ Mesmo \textit{score}!
    \end{block}
  \end{frame}

  %--- SUBSECTION: ALiBi ---
  \subsection{ALiBi}

  %--- Slide: ALiBi Paper ---
  \begin{frame}{ALiBi: Artigo}
    \begin{block}{Referência}
      Press, O. et al. ``Train Short, Test Long: Attention with Linear Biases Enables Input Length Extrapolation.'' ICLR 2022.\footfullcite{press2022alibi}
    \end{block}
    \begin{itemize}
      \item Nenhuma modificação nos \textit{embeddings}: zero vetores posicionais adicionados
      \item Em vez disso: subtrair uma penalidade linear dos \textit{logits} de atenção
      \item Penalidade proporcional a $|i - j|$: \textit{tokens} distantes recebem penalidade maior
      \item Cada \textit{head} $h$ tem uma inclinação diferente $m_h$ $\to$ espectro de atenção local a global
      \item \textcolor{DarkBlue}{\textbf{Extrapolação \textit{zero-shot}: treinado em $L=\text{1K}$, funciona em $L=\text{8K}$}}
    \end{itemize}
  \end{frame}

  %--- Slide: ALiBi Core Idea ---
  \begin{frame}{ALiBi: Ideia Central}
    \textcolor{MidGray}{\small RoPE extrapola apenas com ajustes (interpolação/\textit{fine-tuning}); ALiBi busca extrapolação \textit{zero-shot} por outro caminho, um viés direto na atenção.}

    \vspace{0.2em}
    Atenção padrão (\textit{pre-softmax}):
    \begin{exampleblock}{}
      $A(i, j) = \mathbf{q}_i^\top \mathbf{k}_j \;/\; \sqrt{d}$
    \end{exampleblock}

    Atenção ALiBi (\textit{pre-softmax}):
    \begin{block}{}
      $A_h(i, j) = \mathbf{q}_i^\top \mathbf{k}_j \;/\; \sqrt{d} \;-\; m_h \cdot |i - j|$
    \end{block}

    A penalidade $m_h \cdot |i{-}j|$ cresce com a distância, \textit{tokens} distantes são ponderados para baixo.\\
    Diferentes \textit{heads} $h$ têm inclinações $m_h$ diferentes $\to$ espectro de atenção local a global.

    \begin{itemize}
      \item \textcolor{DarkBlue}{\textbf{Sem vetores de posição, custo de PE é zero}}
      \item \textcolor{DarkBlue}{Em sequências longas de teste, o modelo simplesmente aplica penalidades maiores, permanece \textit{in-distribution}}
    \end{itemize}
  \end{frame}

  %--- Slide: ALiBi Slopes ---
  \begin{frame}{ALiBi: Inclinações por \textit{Head}}
    As inclinações formam uma sequência geométrica sobre os $n$ \textit{attention heads}:

    \begin{block}{}
      $m_h = 2^{-8h/n}$ \qquad $h = 1, 2, \ldots, n$
    \end{block}

    {\small O expoente $8$ é uma escolha empírica de \cite{press2022alibi} (calibra a faixa de inclinações), sem derivação teórica.}

    \centering
    {\small
    \begin{tabular}{ccccc}
      \toprule
      $n$ \textit{heads} & $h{=}1$ (mais íngreme) & $h{=}2$ & $h{=}3$ & $h{=}4$ (mais plano) \\
      \midrule
      4 & $2^{-2} = 0.2500$ & $2^{-4} = 0.0625$ & $2^{-6} = 0.0156$ & $2^{-8} = 0.0039$ \\
      8 & $2^{-1} = 0.5000$ & $2^{-2} = 0.2500$ & $2^{-3} = 0.1250$ & $2^{-4} = 0.0625$ \\
      \bottomrule
    \end{tabular}}

    \flushleft
    \begin{itemize}
      \item Inclinações íngremes ($m$ grande) $\to$ \textit{head} atende principalmente localmente
      \item Inclinações planas ($m$ pequeno) $\to$ \textit{head} atende a longas distâncias
      \item \textcolor{DarkBlue}{Juntos, os \textit{heads} cobrem amplo espectro de raios de atenção}
      \item Inclinações são hiperparâmetros fixos, não aprendidos
    \end{itemize}
  \end{frame}

  %--- Slide: ALiBi Decay Curves Visualization ---
  \begin{frame}{ALiBi: Curvas de Decaimento por \textit{Head}}
    \centering
    \begin{tikzpicture}
      \begin{axis}[
        width=11cm, height=6cm,
        xlabel={distância $|i-j|$},
        ylabel={penalidade $-m_h \cdot |i-j|$},
        xmin=0, xmax=32, ymin=-8, ymax=0.5,
        grid=major, grid style={dashed, gray!30},
        ymajorgrids=true,
        legend style={font=\scriptsize, at={(0.02,0.02)}, anchor=south west, draw=none, fill=white, fill opacity=0.8, text opacity=1},
        label style={font=\footnotesize}, tick label style={font=\scriptsize},
      ]
        % h = 1: steepest
        \addplot[DarkRed, very thick, domain=0:32, samples=2] {-0.25*x};
        \addlegendentry{$h{=}1$, $m{=}0.2500$ (local)}

        % h = 2
        \addplot[DarkBlue, very thick, domain=0:32, samples=2] {-0.0625*x};
        \addlegendentry{$h{=}2$, $m{=}0.0625$}

        % h = 3
        \addplot[MidGray, very thick, domain=0:32, samples=2] {-0.0156*x};
        \addlegendentry{$h{=}3$, $m{=}0.0156$}

        % h = 4: flattest
        \addplot[ForestGreen, very thick, domain=0:32, samples=2] {-0.0039*x};
        \addlegendentry{$h{=}4$, $m{=}0.0039$ (global)}

        \node[font=\scriptsize, DarkRed, anchor=west] at (axis cs:16, -4.5) {\textbf{atenção local}};
        \node[font=\scriptsize, ForestGreen, anchor=west] at (axis cs:16, -0.35) {\textbf{atenção global}};
      \end{axis}
    \end{tikzpicture}

    \vspace{0.3em}
    {\footnotesize
    Inclinações íngremes $\to$ \textit{head} atende principalmente localmente.
    Inclinações planas $\to$ \textit{head} mantém atenção em distâncias longas.\\
    \textcolor{DarkBlue}{\textbf{Juntos, os $n$ \textit{heads} cobrem todo o espectro de raios de atenção, do local ao global.}}
    }
  \end{frame}

  %--- Slide: ALiBi Hand Comp Step 1 ---
  \begin{frame}{Cálculo Manual, ALiBi, Passo 1: Configuração e Inclinações}
    \textbf{Parâmetros:} $n_{\text{heads}} = 4$, \; $\text{seq\_len} = 4$ \; (posições $i, j \in \{0,1,2,3\}$)

    \vspace{0.3em}
    Calcular inclinações (foco nos \textit{heads} 1 e 2):

    \centering
    {\small
    \begin{tabular}{ccc}
      \toprule
      \textit{Head} $h$ & Fórmula $2^{-8h/n}$ & Valor \\
      \midrule
      \rowcolor{CellGreen} 1 & $2^{-8 \cdot 1/4} = 2^{-2}$ & $m_1 = 0.2500$ \\
      \rowcolor{CellGreen} 2 & $2^{-8 \cdot 2/4} = 2^{-4}$ & $m_2 = 0.0625$ \\
      3 & $2^{-8 \cdot 3/4} = 2^{-6}$ & $m_3 = 0.0156$ \\
      4 & $2^{-8 \cdot 4/4} = 2^{-8}$ & $m_4 = 0.0039$ \\
      \bottomrule
    \end{tabular}}
  \end{frame}

  %--- Slide: ALiBi Hand Comp Step 2 ---
  \begin{frame}{Cálculo Manual, ALiBi, Passo 2: Matrizes de \textit{Bias}}
    \begin{columns}
      \begin{column}{0.45\textwidth}
        Matriz $|i - j|$ crua (causal: triângulo superior mascarado):

        {\scriptsize
        \begin{tabular}{c|cccc}
          \toprule
          $i \backslash j$ & $j{=}0$ & $j{=}1$ & $j{=}2$ & $j{=}3$ \\
          \midrule
          $i{=}0$ & 0 & $-\infty$ & $-\infty$ & $-\infty$ \\
          $i{=}1$ & 1 & 0 & $-\infty$ & $-\infty$ \\
          $i{=}2$ & 2 & 1 & 0 & $-\infty$ \\
          $i{=}3$ & 3 & 2 & 1 & 0 \\
          \bottomrule
        \end{tabular}}
      \end{column}
      \begin{column}{0.52\textwidth}
        $\times\; m_1 = 0.25$ $\to$ \textit{Bias} do \textit{head} 1:

        {\scriptsize
        \begin{tabular}{c|cccc}
          \toprule
          $i \backslash j$ & $j{=}0$ & $j{=}1$ & $j{=}2$ & $j{=}3$ \\
          \midrule
          \rowcolor{CellGreen} $i{=}0$ & 0 & $-\infty$ & $-\infty$ & $-\infty$ \\
          \rowcolor{CellGreen} $i{=}1$ & $-0.25$ & 0 & $-\infty$ & $-\infty$ \\
          \rowcolor{CellGreen} $i{=}2$ & $-0.50$ & $-0.25$ & 0 & $-\infty$ \\
          \rowcolor{CellGreen} $i{=}3$ & $-0.75$ & $-0.50$ & $-0.25$ & 0 \\
          \bottomrule
        \end{tabular}}
      \end{column}
    \end{columns}

    \vspace{0.5em}
    Posições próximas ($|i{-}j|$ pequeno) recebem penalidade menor que posições distantes.\\
    \textit{Head} 2 ($m=0.0625$) teria penalidades $4\times$ menores $\to$ raio de atenção efetivo mais longo.
  \end{frame}

  %--- Slide: ALiBi Hand Comp Step 3 ---
  \begin{frame}{Cálculo Manual, ALiBi, Passo 3: Aplicar \textit{Bias} e \textit{Softmax}}
    \textit{Logits} brutos para linha $i{=}2$: \; $[2.0,\; 1.5,\; 3.0,\; -]$

    \vspace{0.3em}
    \resizebox{\linewidth}{!}{%
    \begin{tabular}{ccccc}
      \toprule
      \textit{Head} & \textit{Scores} brutos & \textit{Bias} ALiBi & Após \textit{bias} & $\text{softmax} \approx$ \\
      \midrule
      \rowcolor{CellGreen} $h_1$ ($m{=}0.25$) & $[2.0,\; 1.5,\; 3.0]$ & $[-0.50,\; -0.25,\; 0]$ & $[1.50,\; 1.25,\; 3.0]$ & $[0.16,\; 0.13,\; 0.71]$ \\
      \rowcolor{CellGreen} $h_2$ ($m{=}0.0625$) & $[2.0,\; 1.5,\; 3.0]$ & $[-0.125,\; -0.063,\; 0]$ & $[1.88,\; 1.44,\; 3.0]$ & $[0.21,\; 0.14,\; 0.65]$ \\
      \bottomrule
    \end{tabular}}

    \vspace{0.5em}
    \textit{Head} 1 (inclinação maior): mais peso na posição $j{=}2$ (o \textit{token} imediatamente anterior).\\
    \textit{Head} 2 (inclinação menor): distribui atenção mais amplamente entre $j \in \{0,1,2\}$.

    \vspace{0.3em}
    \textcolor{DarkBlue}{\textbf{Em tempo de teste com $L > L_{\text{train}}$, o \textit{bias} simplesmente fica maior para distâncias maiores $\to$ nenhum problema de OOD.}}
  \end{frame}

  %--- Slide: ALiBi Properties ---
  \begin{frame}{ALiBi: Propriedades}
    \begin{itemize}
      \item Zero custo no \textit{embedding}: nenhum vetor adicionado à entrada
      \item \textcolor{DarkBlue}{\textbf{Extrapolação \textit{zero-shot}: treinado em $L{=}1$K, funciona em $L{=}8$K sem \textit{fine-tuning}}}
      \item Modificação simples: apenas 1--2 linhas de código na atenção
      \item Diferentes \textit{heads} atendem a diferentes faixas automaticamente
      \item \textcolor{DarkRed}{\textbf{Ressalva: decaimento linear pode não ser adequado para todas as tarefas}}
      \item \textcolor{DarkRed}{Inclinações são fixas, não aprendidas por tarefa}
      \item Não é trivialmente compatível com extensão de contexto baseada em RoPE
    \end{itemize}
  \end{frame}

  %--- SUBSECTION: NoPE ---
  \subsection{NoPE}

  %--- Slide: NoPE Paper ---
  \begin{frame}{NoPE: Artigo}
    \begin{block}{Referência}
      Kazemnejad, A. et al. ``The Impact of Positional Encoding on Length Generalization in Transformers.'' NeurIPS 2023.\footfullcite{kazemnejad2023nope}
    \end{block}
    \begin{itemize}
      \item Estudo sistemático: comparar sinusoidal, RoPE, ALiBi e \textbf{nenhum PE}
      \item \textcolor{DarkBlue}{\textbf{Achado surpreendente: modelos sem PE generalizam para sequências mais longas}}
      \item A máscara causal de atenção codifica implicitamente a ordem dos \textit{tokens}
      \item NoPE superou todos os PEs explícitos testados em \textit{benchmarks} de generalização
      \item \textcolor{DarkRed}{\textbf{Ressalvas: resultados dependem fortemente do tipo de tarefa e tamanho do modelo}}
    \end{itemize}
  \end{frame}

  %--- Slide: NoPE Core Idea ---
  \begin{frame}{NoPE: Ideia Central}
    \textcolor{DarkBlue}{\textbf{Não usar nenhum \textit{positional encoding}. Deixar a máscara causal fazer o trabalho.}}

    \begin{itemize}
      \item Em um \textit{decoder} autorregressivo, a máscara de atenção é estritamente triangular inferior
      \item \textit{Token} na posição $i$ só pode atender a posições $j \le i$
      \item \textcolor{DarkBlue}{\textbf{O padrão de quais \textit{tokens} podem ser atendidos é, em si, informação posicional}}
      \item Cada camada vê um ``horizonte de histórico'' diferente
      \item Como a máscara muda suavemente com o comprimento, sequências mais longas não são OOD
    \end{itemize}

    \begin{block}{Por que isso não contradiz o início da aula}
      A invariância a permutação vale para a atenção \textbf{bidirecional}. A \textbf{máscara causal} já quebra essa simetria (cada posição $i$ vê um conjunto diferente de \textit{tokens}), então o \textit{decoder-only} carrega um sinal de ordem mesmo sem PE. Não vale para \textit{encoder-only} (BERT).
    \end{block}
  \end{frame}

  %--- Slide: NoPE Why ---
  \begin{frame}{NoPE: Por que Remover PE Ajuda na Generalização de Comprimento}
    \begin{itemize}
      \item \textcolor{DarkRed}{\textbf{PE absoluto: posições $> L_{\text{train}}$ são \textit{out-of-distribution}}}
      \item \textcolor{DarkRed}{PE relativo (RoPE): frequências calibradas para o comprimento de treino, pode degradar}
      \item \textcolor{DarkBlue}{\textbf{NoPE: nenhum sinal posicional adicionado $\to$ nada se torna OOD em comprimentos maiores}}
      \item O único \textit{inductive bias} posicional é o padrão da máscara, que escala naturalmente
      \item \textcolor{DarkRed}{\textbf{Ressalva: o modelo precisa inferir posição puramente a partir do contexto}}
        \begin{itemize}
          \item Funciona bem para tarefas de linguagem com pistas contextuais ricas
          \item \textcolor{DarkRed}{Falha em tarefas que exigem posição absoluta (indexação, aritmética)}
        \end{itemize}
    \end{itemize}
  \end{frame}

  %--- Slide: NoPE Limitations ---
  \begin{frame}{NoPE: Limitações}
    \begin{itemize}
      \item \textcolor{DarkRed}{\textbf{Depende de máscara causal, não se aplica a modelos \textit{encoder-only} (BERT)}}
      \item \textcolor{DarkRed}{Não lida com tarefas que exigem posição absoluta}
      \item Convergência mais lenta, o modelo precisa aprender padrões posicionais apenas dos dados
      \item Empiricamente pior que RoPE em alguns \textit{benchmarks} de contexto longo (SCROLLS, LongBench)
      \item Métodos de extensão de contexto para RoPE não se aplicam diretamente
      \item \textbf{NoPE é um \textit{baseline} útil e forte generalizador, ainda não é o padrão para grandes LLMs}
    \end{itemize}
  \end{frame}

  %--- Slide: NoPE Discussion ---
  \begin{frame}{NoPE, Discussão: Máscara Causal como PE Implícito}
    \begin{columns}
      \begin{column}{0.45\textwidth}
        Máscara causal para $\text{seq\_len} = 4$:

        {\small
        \begin{tabular}{c|cccc}
          \toprule
          $i \backslash j$ & $j{=}0$ & $j{=}1$ & $j{=}2$ & $j{=}3$ \\
          \midrule
          $i{=}0$ & \cellcolor{CellGreen}1 & \cellcolor{LightGray}--- & \cellcolor{LightGray}--- & \cellcolor{LightGray}--- \\
          $i{=}1$ & \cellcolor{CellGreen}1 & \cellcolor{CellGreen}1 & \cellcolor{LightGray}--- & \cellcolor{LightGray}--- \\
          $i{=}2$ & \cellcolor{CellGreen}1 & \cellcolor{CellGreen}1 & \cellcolor{CellGreen}1 & \cellcolor{LightGray}--- \\
          $i{=}3$ & \cellcolor{CellGreen}1 & \cellcolor{CellGreen}1 & \cellcolor{CellGreen}1 & \cellcolor{CellGreen}1 \\
          \bottomrule
        \end{tabular}}
      \end{column}
      \begin{column}{0.52\textwidth}
        \begin{itemize}
          \item Cada linha tem padrão de máscara único
          \item \textbf{Softmax normaliza sobre $i{+}1$ tokens visíveis:} a magnitude \textbf{média} do peso de atenção escala como $\bar\alpha_{ij} \approx \frac{1}{i{+}1}$ \textcolor{DarkBlue}{$\to$ dá ao modelo um sinal indireto de $i$}
          \item $i{=}0$ vê apenas ele mesmo; $i{=}3$ vê todos os 4
          \item Esse padrão muda previsivelmente com $L$ $\to$ generaliza para $L > L_{\text{train}}$
          \item \textcolor{DarkBlue}{Cada camada empilha esse sinal, modelos profundos ganham representação posicional implícita rica}
        \end{itemize}
      \end{column}
    \end{columns}
  \end{frame}

  %--- Slide: NoPE Multi-Scale Causal Mask Visualization ---
  \begin{frame}{NoPE: Padrão Causal Escalando com $L$}
    \centering
    \begin{tikzpicture}[scale=1.0, every node/.style={font=\tiny}]
      % ============ L = 4 ============
      \begin{scope}[shift={(0,0)}]
        \foreach \i in {0,...,3} {
          \foreach \j in {0,...,3} {
            \pgfmathtruncatemacro{\visible}{\j<=\i ? 1 : 0}
            \ifnum\visible=1
              \fill[CellGreen, draw=gray!60] (\j*0.35, -\i*0.35) rectangle ++(0.35,-0.35);
            \else
              \fill[LightGray, draw=gray!60] (\j*0.35, -\i*0.35) rectangle ++(0.35,-0.35);
            \fi
          }
        }
        \node[font=\scriptsize] at (0.7, -1.65) {$L=4$};
      \end{scope}

      % ============ L = 8 ============
      \begin{scope}[shift={(3.0,0)}]
        \foreach \i in {0,...,7} {
          \foreach \j in {0,...,7} {
            \pgfmathtruncatemacro{\visible}{\j<=\i ? 1 : 0}
            \ifnum\visible=1
              \fill[CellGreen, draw=gray!60] (\j*0.25, -\i*0.25) rectangle ++(0.25,-0.25);
            \else
              \fill[LightGray, draw=gray!60] (\j*0.25, -\i*0.25) rectangle ++(0.25,-0.25);
            \fi
          }
        }
        \node[font=\scriptsize] at (1.0, -2.35) {$L=8$};
      \end{scope}

      % ============ L = 16 ============
      \begin{scope}[shift={(6.5,0)}]
        \foreach \i in {0,...,15} {
          \foreach \j in {0,...,15} {
            \pgfmathtruncatemacro{\visible}{\j<=\i ? 1 : 0}
            \ifnum\visible=1
              \fill[CellGreen, draw=gray!30] (\j*0.16, -\i*0.16) rectangle ++(0.16,-0.16);
            \else
              \fill[LightGray, draw=gray!30] (\j*0.16, -\i*0.16) rectangle ++(0.16,-0.16);
            \fi
          }
        }
        \node[font=\scriptsize] at (1.28, -2.85) {$L=16$};
      \end{scope}
    \end{tikzpicture}

    \vspace{0.8em}
    \begin{itemize}
      \item O padrão triangular é \textbf{auto-similar}: escala continuamente com $L$
      \item \textcolor{DarkBlue}{\textbf{Nenhuma quebra de distribuição em $L > L_{\text{train}}$ $\to$ por isso NoPE generaliza}}
      \item Cada linha $i$ permanece única (vê exatamente $i{+}1$ posições passadas), independente de $L$
    \end{itemize}
  \end{frame}

  %--- SUBSECTION: Comparison ---
  \subsection{Comparação}

  %--- Slide: Comparison Table ---
  \begin{frame}{Comparação: Todos os Métodos de PE}
    \centering
    \resizebox{\linewidth}{!}{%
    \begin{tabular}{lcccc}
      \toprule
      \textbf{Método} & \textbf{Aprendido?} & \textbf{Modifica\ldots} & \textbf{Extrapolação} & \textbf{Custo} \\
      \midrule
      Sinusoidal     & Não & \textit{Input embedding} & Pobre (OOD)           & Nenhum \\
      \textit{Learned Abs.} & Sim & \textit{Input embedding} & Falha (sem novos pos.) & $L \times d$ parâmetros \\
      RoPE           & Não & Rotações Q, K            & Moderada              & Tabela cos/sin \\
      ALiBi          & Não & \textit{Logits} de atenção & Boa (\textit{zero-shot}) & Desprezível \\
      NoPE           & Não & Nada                     & Boa (apenas máscara)  & Nenhum \\
      \bottomrule
    \end{tabular}}

    \vspace{0.5em}
    \flushleft
    \textcolor{DarkBlue}{\textbf{RoPE}} é o padrão \textit{de facto} para grandes LLMs (equilíbrio entre qualidade e extensibilidade).\\
    \textcolor{DarkBlue}{\textbf{ALiBi}} oferece a melhor extrapolação \textit{zero-shot}.\\
    \textcolor{DarkBlue}{\textbf{NoPE}} é um \textit{baseline} forte para generalização de comprimento.
  \end{frame}

  %--- Slide: Attention Patterns ---
  \begin{frame}{Padrões de Peso de Atenção por Método}
    Como cada método modela a atenção em função da distância $|i - j|$:

    \vspace{0.3em}
    \centering
    \resizebox{\linewidth}{!}{%
    \begin{tabular}{lll}
      \toprule
      \textbf{Método} & \textbf{Fórmula do \textit{score}} & \textbf{Efeito em $|i{-}j|$ grande} \\
      \midrule
      Sinusoidal & $\mathbf{q}_i^\top \mathbf{k}_j / \sqrt{d}$ (\textit{rotated embed}) & Imprevisível (absoluto) \\
      RoPE       & $f(\mathbf{q}_i, \mathbf{k}_j, j{-}i)$ via rotação & Decaimento suave (dependente de freq.) \\
      ALiBi      & $\mathbf{q}_i^\top \mathbf{k}_j / \sqrt{d} - m \cdot |i{-}j|$ & Decaimento linear garantido por \textit{head} \\
      NoPE       & $\mathbf{q}_i^\top \mathbf{k}_j / \sqrt{d}$ (sem PE) & Aprendido pelo modelo via contexto \\
      \bottomrule
    \end{tabular}}

    \flushleft
    \begin{itemize}
      \item ALiBi oferece o decaimento mais previsível, mais fácil de raciocinar sobre extrapolação
      \item RoPE tem decaimento dependente de frequência, adaptável via reescalonamento (LongRoPE)
      \item NoPE depende inteiramente do modelo aprender consciência posicional dos dados
    \end{itemize}
  \end{frame}

  %--- Slide: Attention Decay Comparison Visualization ---
  \begin{frame}{Padrões de Decaimento: Comparação Visual}
    \centering
    \begin{tikzpicture}
      \begin{axis}[
        width=11cm, height=6cm,
        xlabel={distância $|i-j|$},
        ylabel={peso de atenção (esquemático)},
        xmin=0, xmax=50, ymin=0, ymax=1.05,
        grid=major, grid style={dashed, gray!30},
        ymajorgrids=true,
        legend style={font=\scriptsize, at={(0.98,0.98)}, anchor=north east, draw=none, fill=white, fill opacity=0.8, text opacity=1},
        label style={font=\footnotesize}, tick label style={font=\scriptsize},
      ]
        % OOD shaded region
        \fill[LightGray, opacity=0.6] (axis cs:30,0) rectangle (axis cs:50,1.05);
        \node[font=\scriptsize, MidGray, align=center] at (axis cs:40,0.95) {região de\\extrapolação};
        \draw[dashed, MidGray] (axis cs:30,0) -- (axis cs:30,1.05);

        % Sinusoidal: oscillatory, imprevisible
        \addplot[MidGray, thick, dashed, domain=0:50, samples=120]
          {0.55 + 0.35*cos(deg(0.8*x)) * exp(-0.02*x)};
        \addlegendentry{Sinusoidal (imprevisível)}

        % RoPE: smooth exponential decay
        \addplot[DarkBlue, very thick, domain=0:50, samples=80]
          {exp(-0.06*x)};
        \addlegendentry{RoPE (suave)}

        % ALiBi: linear decay (clipped at 0)
        \addplot[DarkRed, very thick, domain=0:50, samples=60]
          {max(0, 1 - 0.025*x)};
        \addlegendentry{ALiBi (linear)}

        % NoPE: learned shape, bumpy
        \addplot[ForestGreen, very thick, dotted, domain=0:50, samples=120]
          {0.9*exp(-0.04*x) + 0.08*sin(deg(0.4*x))};
        \addlegendentry{NoPE (aprendido)}
      \end{axis}
    \end{tikzpicture}

    \vspace{0.2em}
    {\footnotesize
    \textcolor{DarkRed}{\textbf{ALiBi}}: decaimento mais previsível $\to$ extrapolação garantida.\quad
    \textcolor{DarkBlue}{\textbf{RoPE}}: decaimento suave mas dependente de frequência.\\
    Sinusoidal tende a oscilar; NoPE depende do que o modelo aprendeu.\\[0.2em]
    \textit{\tiny Curvas esquemáticas para comparação qualitativa, não valores experimentais.}
    }
  \end{frame}

  %--- Slide: Empirical Results ---
  \begin{frame}{Resultados Empíricos: Perplexidade vs.\ Comprimento de Contexto}
    {\small Desempenho relativo em modelagem de linguagem (menor = melhor). Todos treinados em $L = 2048$:}

    \centering
    \begin{tikzpicture}
      \begin{axis}[
        xlabel={Comprimento de contexto},
        ylabel={Perplexidade},
        symbolic x coords={512,1K,2K,3K,4K,6K,8K},
        xtick=data,
        ymajorgrids=true,
        grid style=dashed,
        width=0.85\linewidth,
        height=5.5cm,
        legend style={at={(0.23,0.98)}, anchor=north east, font=\scriptsize},
        thick,
      ]
        \addplot[color=MidGray, mark=square*] coordinates {(512,12)(1K,12.5)(2K,13)(3K,22)(4K,60)(6K,150)(8K,400)};
        \addplot[color=DarkBlue, mark=*] coordinates {(512,12)(1K,12.4)(2K,12.8)(3K,14)(4K,17)(6K,28)(8K,55)};
        \addplot[color=DarkRed, mark=triangle*] coordinates {(512,12)(1K,12.4)(2K,12.9)(3K,13.2)(4K,13.5)(6K,14)(8K,14.5)};
        \addplot[color=ForestGreen, mark=diamond*] coordinates {(512,12)(1K,12.4)(2K,12.9)(3K,13.1)(4K,13.4)(6K,13.7)(8K,14)};
        \legend{Sinusoidal, RoPE, ALiBi, NoPE}
      \end{axis}
    \end{tikzpicture}

    {\scriptsize\textit{Ilustrativo; tendências relativas de Press et al.\ (2022), Kazemnejad et al.\ (2023)}}
  \end{frame}

  %--- Slide: Transition ---
  \begin{frame}{Transição: Do Design de PE para Desafios de Uso}
    Agora temos modelos com bons \textit{positional encodings}.\\
    Mas duas questões-chave permanecem:

    \vspace{0.5em}
    \begin{columns}[t]
      \begin{column}{0.45\textwidth}
        \begin{alertblock}{1. Contexto longo realmente ajuda?}
          \textit{Lost in the Middle}\\
          $\to$ Parte 3
        \end{alertblock}
      \end{column}
      \begin{column}{0.45\textwidth}
        \begin{block}{2. Podemos estender contexto além do treino?}
          LongRoPE, YaRN, PI\\
          $\to$ Parte 4
        \end{block}
      \end{column}
    \end{columns}
  \end{frame}

  %======================================================================================
  % SECTION 3: LONG-CONTEXT CHALLENGES
  %======================================================================================
  \section{Desafios em Contextos Longos}
  \subsection{Perplexidade e suas Limitações em Contexto Longo}

  %--- Slide: Perplexity definition ---
  \begin{frame}{Perplexidade: Definição}
    \textbf{Definição}: mede o quão ``surpreso'' o modelo fica ao ver o texto.
    \[
      \mathrm{PPL}(w_1,\ldots,w_N)
        = \exp\!\left(-\frac{1}{N}\sum_{i=1}^{N}\log P(w_i \mid w_1,\ldots,w_{i-1})\right)
    \]

    \vspace{0.3em}
    \begin{columns}[T]
      \begin{column}{0.48\textwidth}
        \begin{block}{Vantagens}
          \begin{itemize}
            \item Simples de calcular em qualquer corpus
            \item Permite comparar PEs em condições iguais
            \item Correlaciona com qualidade do modelo de linguagem
          \end{itemize}
        \end{block}
      \end{column}
      \begin{column}{0.48\textwidth}
        \begin{alertblock}{Atenção}
          PPL é uma \textbf{média sobre todos os tokens}.\\[0.3em]
          Tokens iniciais são previstos com contexto curto; tokens finais com contexto completo.\\[0.3em]
          Contribuições locais e distantes ficam \textit{misturadas}.
        \end{alertblock}
      \end{column}
    \end{columns}
  \end{frame}

  %--- Slide: Why perplexity is not enough for long context ---
  \begin{frame}{Perplexidade Não Captura Uso Real do Contexto Longo}
    \begin{alertblock}{O problema central}
      Um modelo com \textbf{PPL baixa} pode \textbf{ignorar completamente} o contexto distante, e ainda assim parecer ``bom''.
    \end{alertblock}

    \vspace{0.3em}
    \textbf{Por quê?}
    \begin{itemize}
      \item A maioria dos tokens é prevista corretamente a partir de contexto \textbf{local} (bigramas, padrões sintáticos)
      \item PPL \textit{dilui} falhas de longo alcance na média global
      \item Um modelo com PPL = 14 a 8K tokens pode \textbf{nunca usar} o que foi dito nos primeiros 6K
    \end{itemize}

    \vspace{0.3em}
    \begin{columns}[T]
      \begin{column}{0.48\textwidth}
        \begin{block}{O que PPL mede}
          Qualidade de \textit{previsão do próximo token} em média sobre todos os tokens
        \end{block}
      \end{column}
      \begin{column}{0.48\textwidth}
        \begin{alertblock}{O que PPL \textit{não} mede}
          Se o modelo \textit{utiliza} informações de posições distantes ao gerar uma resposta
        \end{alertblock}
      \end{column}
    \end{columns}

    \vspace{0.3em}
    \textcolor{DarkRed}{$\Rightarrow$ Precisamos de avaliações baseadas em tarefas de recuperação de longo alcance.}
  \end{frame}

  \subsection{Lost in the Middle}

  %--- Slide: Lost Paper ---
  \begin{frame}{\textit{Lost in the Middle}: Artigo}
    \begin{block}{Referência}
      Liu, N. F. et al. ``Lost in the Middle: How Language Models Use Long Contexts.'' TACL 2024.\footfullcite{liu2024lost}
    \end{block}
    \begin{itemize}
      \item Pergunta: onde colocar a informação relevante em um contexto longo?
      \item Tarefa: QA multi-documento, um documento contém a resposta
      \item Variou: (a) número de documentos, (b) posição do documento relevante
      \item Testou GPT-3.5-Turbo, GPT-4, Claude e modelos \textit{open-source}
      \item \textcolor{DarkRed}{\textbf{Todos os modelos exibem curva de desempenho em forma de U sobre a posição}}
    \end{itemize}
  \end{frame}

  %--- Slide: U-curve ---
  \begin{frame}{\textit{Lost in the Middle}: Curva de Desempenho em U}
    Acurácia em QA multi-documento em função da posição do documento relevante:

    \centering
    \begin{tikzpicture}
      \begin{axis}[
        xlabel={Posição do documento relevante (\%)},
        ylabel={Acurácia (\%)},
        symbolic x coords={10\%,20\%,30\%,40\%,50\%,60\%,70\%,80\%,90\%},
        xtick=data,
        ymajorgrids=true,
        grid style=dashed,
        width=0.85\linewidth,
        height=5.5cm,
        mark=*,
        thick,
        color=DarkBlue,
        ymin=30, ymax=80,
      ]
        \addplot coordinates {(10\%,72)(20\%,63)(30\%,52)(40\%,45)(50\%,43)(60\%,46)(70\%,53)(80\%,62)(90\%,71)};
      \end{axis}
    \end{tikzpicture}

    Pico no início e fim (\textit{primacy / recency}), \textcolor{DarkRed}{\textbf{vale no meio}}.

    {\scriptsize\textit{Redesenhado a partir de Liu et al., Lost in the Middle (TACL 2024)}}
  \end{frame}

  %--- Slide: Lost Why ---
  \begin{frame}{\textit{Lost in the Middle}: Por que Isso Acontece?}
    \begin{itemize}
      \item \textbf{Efeito de primazia}: \textit{tokens} próximos à posição 0 recebem alta atenção nas camadas iniciais
        \begin{itemize}
          \item A máscara causal força todos os \textit{tokens} a atender ao início da sequência
        \end{itemize}
      \item \textbf{Efeito de recência}: os últimos \textit{tokens} estão na ``memória de trabalho'' ativa da atenção
      \item \textcolor{DarkRed}{\textbf{\textit{Tokens} do meio, grande distância $|i{-}j|$ de ambas as extremidades}}
        \begin{itemize}
          \item Mesmo sem decaimento explícito de PE, padrões de atenção aprendidos tendem a ser locais
        \end{itemize}
      \item \textcolor{DarkRed}{\textbf{O alcance efetivo de atenção do modelo é menor que sua janela de contexto}}
      \item \textcolor{DarkBlue}{\textbf{Nota:} GPT-4, Claude e LLaMA-2 usam RoPE ou variantes, o fenômeno \textbf{não} é resolvido trocando o PE. É um padrão aprendido de atenção, não uma propriedade do \textit{encoding}.}
    \end{itemize}
  \end{frame}

  %--- Slide: Lost Implications ---
  \begin{frame}{\textit{Lost in the Middle}: Implicações Práticas}
    \begin{itemize}
      \item \textcolor{DarkRed}{\textbf{Simplesmente aumentar o comprimento do contexto NÃO é suficiente}}
        \begin{itemize}
          \item Um modelo com contexto 32K pode ter desempenho pior que um de 4K se a resposta estiver no meio
        \end{itemize}
      \item \textbf{A estratégia de posicionamento dos documentos importa enormemente}
        \begin{itemize}
          \item \textcolor{DarkBlue}{Colocar o contexto mais importante no início ou no fim}
          \item \textcolor{DarkBlue}{RAG (\textit{Retrieval-Augmented Generation}): buscar apenas o trecho relevante}
        \end{itemize}
      \item \textbf{Extensão de contexto (Parte 4) precisa ser combinada com dados de treino de contexto longo}
    \end{itemize}
  \end{frame}

  %--- SUBSECTION: Needle in a Haystack ---
  \subsection{Needle in a Haystack}

  %--- Slide: NiH Setup ---
  \begin{frame}{\textit{Needle in a Haystack}: Configuração}
    \begin{exampleblock}{Referência}
      Kamradt, G. ``LLMTest\_NeedleInAHaystack'' (2023), \textit{benchmark} da comunidade\footfullcite{kamradt2023needle}
    \end{exampleblock}

    \begin{itemize}
      \item Teste de estresse sintético para recuperação em contextos longos
      \item \textcolor{DarkBlue}{\textbf{\textit{Needle}:}} fato curto e específico inserido em posição controlada
        \begin{itemize}
          \item Ex: ``\texttt{A melhor pizzaria de Redmond é a Giovanni's.}''
        \end{itemize}
      \item \textbf{\textit{Haystack}:} texto irrelevante e longo preenchendo o contexto
      \item \textbf{\textit{Depth} \%:} fração do contexto antes da \textit{needle} (0\% = início, 100\% = fim)
      \item \textbf{Varredura:} comprimento do contexto (eixo x) $\times$ \textit{depth} \% (eixo y) $\to$ \textit{heatmap} 2D
    \end{itemize}
  \end{frame}

  %--- Slide: NiH Heatmap ---
  \begin{frame}{\textit{Needle in a Haystack}: Como Ler o \textit{Heatmap}}
    \begin{columns}
      \begin{column}{0.48\textwidth}
        \textbf{Eixos do \textit{heatmap}:}
        \begin{itemize}
          \item Eixo x: comprimento do contexto (1K $\to$ 128K)
          \item Eixo y: \textit{depth} \% da \textit{needle}
          \item Cor: acurácia de recuperação
        \end{itemize}

        \vspace{0.3em}
        \textbf{Padrões de falha:}
        \begin{itemize}
          \item \textcolor{DarkRed}{Faixa vermelha em contexto longo + meio $\to$ \textit{Lost-in-the-Middle}}
          \item \textcolor{DarkRed}{Vermelho em todas as profundidades $\to$ limite rígido de contexto}
          \item \textcolor{DarkBlue}{Verde em todo lugar $\to$ modelo ideal}
        \end{itemize}
      \end{column}
      \begin{column}{0.48\textwidth}
        {\scriptsize Padrão ilustrativo (modelo fraco):}

        \vspace{0.2em}
        \resizebox{\linewidth}{!}{%
        \begin{tabular}{c|cccc}
          \toprule
          \textit{depth} $\backslash$ comp. & 4K & 16K & 64K & 128K \\
          \midrule
          10\% (início)   & \cellcolor{CellGreen}$\checkmark$ & \cellcolor{CellGreen}$\checkmark$ & \cellcolor{CellGreen}$\checkmark$ & \cellcolor{CellGreen}$\checkmark$ \\
          30\%             & \cellcolor{CellGreen}$\checkmark$ & \cellcolor{CellGreen}$\checkmark$ & \cellcolor{CellYellow}$\sim$ & \cellcolor{CellRed}$\times$ \\
          50\% (meio)      & \cellcolor{CellGreen}$\checkmark$ & \cellcolor{CellYellow}$\sim$ & \cellcolor{CellRed}$\times$ & \cellcolor{CellRed}$\times$ \\
          70\%             & \cellcolor{CellGreen}$\checkmark$ & \cellcolor{CellGreen}$\checkmark$ & \cellcolor{CellYellow}$\sim$ & \cellcolor{CellRed}$\times$ \\
          90\% (fim)       & \cellcolor{CellGreen}$\checkmark$ & \cellcolor{CellGreen}$\checkmark$ & \cellcolor{CellGreen}$\checkmark$ & \cellcolor{CellYellow}$\sim$ \\
          \bottomrule
        \end{tabular}}
      \end{column}
    \end{columns}
  \end{frame}

  %--- Slide: NiH Results ---
  \begin{frame}{\textit{Needle in a Haystack}: Resultados Típicos}
    \centering
    \resizebox{\linewidth}{!}{%
    \begin{tabular}{lccc}
      \toprule
      \textbf{Modelo / PE} & \textbf{Contexto} & \textbf{Meio a 64K} & \textbf{Padrão geral} \\
      \midrule
      GPT-4 / RoPE     & 128K   & $\approx$ 90--100\%  & Quase perfeito \\
      GPT-3.5 / RoPE   & 16K    & 50--70\%             & Falha em longo+meio \\
      LLaMA-2 / RoPE   & 4K     & $<$ 20\% a 16K+      & Colapsa além de $L$ \\
      Claude-2 / RoPE  & 100K   & $\approx$ 80\%       & Maioria verde, lacunas no meio \\
      Modelos fracos    & varia  & $<$ 50\% meio        & Zona de falha em U \\
      \bottomrule
    \end{tabular}}

    \flushleft
    \begin{itemize}
      \item Modelos melhores (maiores, mais dados de treino de contexto longo) têm acurácia mais forte no meio
      \item Extensão de contexto (Parte 4) primariamente estende o eixo x (maior $L$)
      \item \textcolor{DarkRed}{\textbf{Estender $L$ sem corrigir o formato U não é suficiente para tarefas reais de contexto longo}}
    \end{itemize}

    {\scriptsize\textit{Valores aproximados de avaliações publicadas de Needle-in-a-Haystack}}
  \end{frame}

  %--- Slide: Takeaway ---
  \begin{frame}{Conclusão: Janela de Contexto $\neq$ Compreensão}
    Uma janela de contexto mais ampla é \textbf{necessária} mas \textcolor{DarkRed}{\textbf{não suficiente}} para bom raciocínio em contextos longos.

    \vspace{0.5em}
    Para realmente se beneficiar de uma janela de contexto longa, um modelo precisa de:
    \begin{itemize}
      \item \textcolor{DarkBlue}{Dados de treino com dependências de longo alcance (não apenas sequências longas)}
      \item \textcolor{DarkBlue}{Padrões de atenção que recuperem informação do meio do contexto de forma confiável}
      \item \textcolor{DarkBlue}{Possivelmente: estratégias de recuperação ou \textit{chunking} para trazer o conteúdo certo}
    \end{itemize}

    \vspace{0.3em}
    \textcolor{DarkRed}{\textbf{Métodos de extensão de contexto estendem a janela, não corrigem o formato U.}}

    O campo está ativamente combinando extensão com melhor treino de contexto longo.
  \end{frame}

  %======================================================================================
  % SECTION 4: CONTEXT EXTENSION METHODS
  %======================================================================================
  \section{Métodos de Extensão de Contexto}

  %--- Slide: Technical View ---
  \begin{frame}{O Problema da Extrapolação: Visão Técnica}
    \textcolor{DarkRed}{\textbf{Por que RoPE falha além de $L_{\text{train}}$?}}

    \begin{itemize}
      \item As frequências $\theta_k$ do RoPE são calibradas para que em $L_{\text{train}}$ todos os ângulos de rotação sejam visitados
      \item Dims de alta frequência (\textit{low-k}): ângulo $i \cdot \theta_k$ dá muitas voltas mesmo para $L_{\text{train}}$
      \item Dims de baixa frequência (\textit{high-k}): ângulo $i \cdot \theta_k$ quase não se move
      \item \textcolor{DarkRed}{\textbf{Além de $L_{\text{train}}$: dims de baixa freq. atingem ângulos NUNCA vistos $\to$ OOD}}
      \item \textcolor{DarkBlue}{\textbf{Solução: reescalar índices de posição para manter o ângulo efetivo \textit{in-distribution}}}
    \end{itemize}

    \begin{exampleblock}{Ângulo efetivo}
      $\varphi(i, k) = i \cdot \theta_k$ \qquad Problema OOD: $\varphi(i, k) > \max$ ângulo de treino quando $i > L_{\text{train}}$
    \end{exampleblock}
  \end{frame}

  %--- Slide: Position Interpolation ---
  \begin{frame}{\textit{Position Interpolation} (PI)}
    \begin{exampleblock}{Referência}
      Chen, S. et al. ``Extending Context Window of LLMs via Positional Interpolation.'' arXiv 2023.\footfullcite{chen2023extending}
    \end{exampleblock}

    Correção mais simples: escalar os índices de posição para ficarem dentro de $[0, L_{\text{train}}]$.

    \begin{block}{Reescalonamento uniforme}
      $i \;\leftarrow\; i \cdot (L_{\text{train}} / L_{\text{test}})$
    \end{block}

    \begin{itemize}
      \item Ângulo efetivo: $\varphi'(i,k) = i \cdot (L_{\text{train}}/L_{\text{test}}) \cdot \theta_k$ $\to$ permanece na faixa de treino
      \item Requer curto \textit{fine-tuning} ($\sim$1000 passos) \textcolor{MidGray}{\scriptsize (adaptar aos ângulos interpolados, OOD no treino)}
      \item \textcolor{DarkBlue}{Alcança contexto de 8K--32K a partir de modelo treinado em 2K}
      \item \textcolor{DarkRed}{\textbf{Problema:} escalonamento uniforme comprime dims de alta freq. desnecessariamente (perde resolução fina)}
    \end{itemize}
  \end{frame}

  %--- Slide: YaRN ---
  \begin{frame}{YaRN: \textit{Yet Another RoPE extensioN}}
    \begin{exampleblock}{Referência}
      Peng, B. et al. ``YaRN: Efficient Context Window Extension of Large Language Models.'' ICLR 2024.\footfullcite{peng2024yarn}
    \end{exampleblock}

    \begin{itemize}
      \item \textbf{Diferentes frequências do RoPE precisam de fatores de escala diferentes}
      \item Dims de alta freq.: já bem cobertas, manter como estão
      \item Dims de baixa freq.: muito esparsas, precisam de interpolação (como PI)
      \item \textcolor{DarkBlue}{Dims intermediárias: escalonamento parcial (\textit{NTK-aware}, do \textit{Neural Tangent Kernel})}
      \item \textcolor{DarkBlue}{\textbf{YaRN aplica escalonamento não-uniforme: escala diferente por dimensão}}
      \item Corrige temperatura de atenção ($\sqrt{d} \to \sqrt{d} \cdot t$) para evitar supressão de \textit{scores}
      \item \textcolor{DarkBlue}{\textbf{Alcança 128K a partir de 2K/4K com $\sim$400 passos de \textit{fine-tuning}; melhor retenção de contexto curto que PI}}
    \end{itemize}
  \end{frame}

  %--- Slide: LongRoPE Paper ---
  \begin{frame}{LongRoPE: Artigo}
    \begin{block}{Referência}
      Ding, Y. et al. ``LongRoPE: Extending LLM Context Window Beyond 2 Million Tokens.'' arXiv 2024.\footfullcite{ding2024longrope}
    \end{block}
    \begin{itemize}
      \item Alvo: estender contexto para 2\,048K \textit{tokens} (2M+) usando LLaMA-2 7B/13B
      \item Inovação: fatores de reescalonamento não-uniformes por dimensão via busca evolucionária
      \item \textit{Fine-tuning} em dois estágios para alcançar 2M preservando qualidade em contexto curto
      \item \textcolor{DarkBlue}{\textbf{Em 2M \textit{tokens}: ainda competitivo em \textit{Needle-in-a-Haystack}}}
      \item Mostra que RoPE é extensível muito além de 100K com o reescalonamento correto
    \end{itemize}
  \end{frame}

  %--- Slide: LongRoPE Uniform Problem ---
  \begin{frame}{LongRoPE: Problema com Escalonamento Uniforme (PI)}
    \textcolor{DarkRed}{\textbf{Por que reescalonamento uniforme (PI) tem desempenho inferior em contextos muito longos?}}

    \vspace{0.3em}
    \centering
    {\small
    \begin{tabular}{cccc}
      \toprule
      \textbf{Índice dim $k$} & $\theta_k$ (freq.) & \textbf{Problema c/ escala uniforme} & \textbf{Tratamento ideal} \\
      \midrule
      Low $k$ (0--15)  & $\approx 1$ (rápida) & Dá muitas voltas $\to$ OK      & Manter inalterado ($s{=}1$) \\
      Mid $k$ (16--55) & $\approx 0.1$        & Parcial $\to$ OOD leve          & Escala pequena \\
      High $k$ (56--63)& $\approx 0.01$ (lenta) & OOD, problema principal    & Escala grande \\
      \bottomrule
    \end{tabular}}

    \smallskip
    PI aplica a mesma escala a \textbf{todas} as dimensões.
    \textcolor{DarkRed}{Sobre-comprime dims \textit{low-k} e sub-comprime \textit{mid-k}.}

    \begin{block}{LongRoPE: reescalonamento por dimensão}
      $i \cdot \theta_k \;\leftarrow\; i \cdot \theta_k / \lambda_k$ \qquad $\lambda_k$ específico por dimensão (busca evolucionária)
    \end{block}
  \end{frame}

  %--- Slide: LongRoPE Evolutionary Search ---
  \begin{frame}{LongRoPE: Reescalonamento Não-Uniforme via Busca Evolucionária}
    \begin{itemize}
      \item \textbf{Objetivo: encontrar fatores de escala por dimensão} $\boldsymbol{\lambda} = [\lambda_0, \lambda_1, \ldots, \lambda_{d/2-1}]$
      \item Função objetivo: minimizar perplexidade em $L_{\text{test}}$ mantendo $\lambda$ em faixa viável
      \item \textcolor{DarkBlue}{\textbf{Método: busca evolucionária (\textbf{CMA-ES}, \textit{Covariance Matrix Adaptation Evolution Strategy}) sobre o vetor $\boldsymbol{\lambda}$}}
      \item População de vetores $\boldsymbol{\lambda}$ candidatos avaliados por métrica de perplexidade
      \item Converge em poucas centenas de gerações, viável na prática
      \item \textcolor{DarkBlue}{\textbf{Resultado: dims \textit{low-k} recebem $\lambda \approx 1$ (inalteradas); dims \textit{high-k} recebem $\lambda$ grande}}
      \item O vetor $\boldsymbol{\lambda}$ final é fixo na inferência, zero custo em tempo de execução
    \end{itemize}
  \end{frame}

  %--- Slide: LongRoPE Angle Trajectories Visualization ---
  \begin{frame}{LongRoPE: Original vs PI vs LongRoPE}
    \centering
    \resizebox{\linewidth}{!}{%
    \begin{tikzpicture}
      \begin{axis}[
        name=origAx,
        width=5.2cm, height=4.2cm,
        title={\textbf{Original}},
        xlabel={posição $i$}, ylabel={$\varphi(i,k) = i \cdot \theta_k$},
        xmin=0, xmax=2000, ymin=0, ymax=2000,
        grid=major, grid style={dashed, gray!30},
        legend style={font=\tiny, at={(0.02,0.98)}, anchor=north west, draw=none, fill=none},
        label style={font=\scriptsize}, tick label style={font=\tiny},
        title style={font=\scriptsize},
      ]
        \fill[LightGray, opacity=0.7] (axis cs:1024,0) rectangle (axis cs:2000,2000);
        \node[font=\tiny, MidGray] at (axis cs:1500,1800) {OOD};
        \draw[dashed, MidGray, thick] (axis cs:1024,0) -- (axis cs:1024,2000);
        \addplot[DarkBlue, very thick, domain=0:2000, samples=60] {x*1.0};
        \addplot[MidGray, very thick, domain=0:2000, samples=60] {x*0.1};
        \addplot[DarkRed, very thick, domain=0:2000, samples=60] {x*0.01};
        \legend{low-$k$, mid-$k$, high-$k$}
      \end{axis}

      \begin{axis}[
        name=piAx,
        at={($(origAx.east)+(1.2cm,0)$)}, anchor=west,
        width=5.2cm, height=4.2cm,
        title={\textbf{PI (uniforme, $s{=}2$)}},
        xlabel={posição $i$},
        xmin=0, xmax=2000, ymin=0, ymax=2000,
        grid=major, grid style={dashed, gray!30},
        label style={font=\scriptsize}, tick label style={font=\tiny},
        title style={font=\scriptsize},
      ]
        \draw[dashed, MidGray, thick] (axis cs:1024,0) -- (axis cs:1024,2000);
        \addplot[DarkBlue, very thick, domain=0:2000, samples=60] {x*1.0/2};
        \addplot[MidGray, very thick, domain=0:2000, samples=60] {x*0.1/2};
        \addplot[DarkRed, very thick, domain=0:2000, samples=60] {x*0.01/2};
        \node[font=\tiny, DarkRed, align=center] at (axis cs:1200,1550) {low-$k$\\sub-utilizada};
      \end{axis}

      \begin{axis}[
        at={($(piAx.east)+(1.2cm,0)$)}, anchor=west,
        width=5.2cm, height=4.2cm,
        title={\textbf{LongRoPE ($\lambda_k$ por dim)}},
        xlabel={posição $i$},
        xmin=0, xmax=2000, ymin=0, ymax=2000,
        grid=major, grid style={dashed, gray!30},
        label style={font=\scriptsize}, tick label style={font=\tiny},
        title style={font=\scriptsize},
      ]
        \draw[dashed, MidGray, thick] (axis cs:1024,0) -- (axis cs:1024,2000);
        \addplot[DarkBlue, very thick, domain=0:2000, samples=60] {x*1.0/1.0};
        \addplot[MidGray, very thick, domain=0:2000, samples=60] {x*0.1/1.8};
        \addplot[DarkRed, very thick, domain=0:2000, samples=60] {x*0.01/6.0};
        \node[font=\tiny, DarkBlue, align=center] at (axis cs:1200,1700) {low-$k$\\intacta};
      \end{axis}
    \end{tikzpicture}}

    \vspace{0.3em}
    {\small
    \textcolor{DarkRed}{\textbf{PI}} trata todas as dimensões igualmente $\to$ compromete \textit{low-k}.\\
    \textcolor{DarkBlue}{\textbf{LongRoPE}} aplica $\lambda_k$ por dimensão $\to$ cada frequência recebe o tratamento certo.
    }
  \end{frame}

  %--- Slide: LongRoPE Two-Stage ---
  \begin{frame}{LongRoPE: \textit{Fine-Tuning} em Dois Estágios}
    \textcolor{DarkRed}{\textbf{Desafio:} \textit{fine-tuning} direto em 2M \textit{tokens} é computacionalmente proibitivo.}

    \vspace{0.3em}
    \centering
    {\small
    \begin{tabular}{cccc}
      \toprule
      \textbf{Estágio} & \textbf{Contexto} & \textbf{Fatores $\lambda$} & \textbf{Descrição} \\
      \midrule
      Base    & 4K (LLaMA-2)  & $\theta_k$ (inalterado)      & Ponto de partida pré-treinado \\
      Estágio 1 & 256K        & $\lambda^{(1)}$ da busca     & FT completo com dados de 256K \\
      Estágio 2 & 2048K       & $\lambda^{(2)}$ da busca     & FT curto em 2M; ajustar $\lambda$ \\
      \bottomrule
    \end{tabular}}

    \flushleft
    \begin{itemize}
      \item Estágio 1 (256K): dados de \textit{fine-tuning} completos; $\lambda^{(1)}$ otimizado para 256K
      \item Estágio 2 (2M): dados mínimos; $\lambda^{(2)}$ re-otimizado para alvo de 2M
      \item \textcolor{DarkBlue}{\textbf{Recuperação de contexto curto: $\lambda$ separado para contextos curtos após estágio 2}}
      \item \textcolor{DarkBlue}{Valores $\lambda$ diferentes para inferência curta vs.\ longa $\to$ chaveamento dinâmico}
    \end{itemize}

    {\scriptsize\textit{Adaptado de Ding et al., LongRoPE (arXiv 2024)}}
  \end{frame}

  %--- Slide: LongRoPE Results ---
  \begin{frame}{LongRoPE: Resultados}
    \centering
    \resizebox{\linewidth}{!}{%
    \begin{tabular}{lcccc}
      \toprule
      \textbf{Método} & \textbf{Ctx máx.} & \textbf{PPL em 2M} & \textbf{PPL 4K} & \textbf{NiH 1M meio} \\
      \midrule
      LLaMA-2 original & 4K    & $\gg$100 (falha)    & 5.5 & $<$ 10\% \\
      PI               & 32K   & $>$50 (degrada)     & 5.6 & 30--50\% \\
      YaRN             & 128K  & $\sim$18            & 5.7 & 65--75\% \\
      LongRoPE         & 2048K & $\sim$9             & 5.6 & $\approx$ 90\% \\
      \bottomrule
    \end{tabular}}

    \flushleft
    \begin{itemize}
      \item \textcolor{DarkBlue}{\textbf{LongRoPE alcança 2M com qualidade comparável em contexto curto}}
      \item \textcolor{DarkBlue}{\textbf{\textit{Needle-in-a-Haystack}: 90\%+ de acurácia em 1M, \textit{depth} 50\%}}
      \item Abordagem em dois estágios adiciona custo modesto de \textit{fine-tuning} vs.\ treino completo em 2M
    \end{itemize}

    {\scriptsize\textit{Dados de Ding et al., LongRoPE (arXiv 2024)}}
  \end{frame}

  %--- Slide: Comparison ---
  \begin{frame}{Métodos de Extensão de Contexto: Comparação}
    \centering
    \resizebox{\linewidth}{!}{%
    \begin{tabular}{lccccc}
      \toprule
      \textbf{Método} & \textbf{Base PE} & \textbf{Tipo de escala} & \textbf{Ctx máx.} & \textbf{Custo de FT} & \textbf{Qual.\ ctx curto} \\
      \midrule
      PI       & RoPE & Uniforme             & 32K--64K & Baixo ($\sim$1K passos)  & Moderada \\
      YaRN     & RoPE & Não-uniforme (NTK)   & 128K     & Baixo ($\sim$400 passos) & Boa \\
      LongRoPE & RoPE & Por dim (evoluída)   & 2M+      & Médio (2 estágios)       & Excelente \\
      \bottomrule
    \end{tabular}}

    \flushleft
    \begin{itemize}
      \item Todos os métodos fazem \textit{fine-tune} de modelo RoPE existente, não treinam do zero
      \item \textit{Trade-off}: extensão mais agressiva $\to$ mais \textit{fine-tuning} e reescalonamento cuidadoso
      \item \textcolor{DarkBlue}{\textbf{Na prática: YaRN é o ponto ideal para a maioria dos casos (128K a baixo custo)}}
      \item Questão aberta: combinar extensão com melhores dados de treino de contexto longo
    \end{itemize}

    {\scriptsize\textit{Compilado de Chen et al.\ (2023), Peng et al.\ (2024), Ding et al.\ (2024)}}
  \end{frame}

  %======================================================================================
  % SUMMARY & REFERENCES
  %======================================================================================
  \section*{Resumo}

  %--- Slide: Summary ---
  \begin{frame}{Resumo}
    \centering
    \resizebox{\linewidth}{!}{%
    \begin{tabular}{llllc}
      \toprule
      \textbf{Método} & \textbf{Tipo} & \textbf{Ideia central} & \textbf{Extrapolação} & \textbf{Custo} \\
      \midrule
      PE Sinusoidal        & Absoluto     & Frequências sin/cos fixas         & Pobre          & Nenhum \\
      \textit{Learned} PE  & Absoluto     & Matriz de \textit{embeddings} treinável & Nenhuma   & $L \!\times\! d$ parâm. \\
      RoPE                 & Relativo     & Rotacionar Q, K por $i \cdot \theta_k$ & Moderada  & Tabela cos/sin \\
      ALiBi                & Relativo     & \textit{Bias} linear $-m \!\cdot\! |i{-}j|$ & Boa      & Desprezível \\
      NoPE                 & Rel./Implíc. & Apenas máscara causal             & Boa            & Nenhum \\
      PI                   & Extensão     & Escalonamento uniforme            & 32K--64K       & FT baixo \\
      YaRN                 & Extensão     & Escala não-uniforme NTK           & 128K           & FT baixo \\
      LongRoPE             & Extensão     & Escala por dim evoluída           & 2M+            & FT médio \\
      \bottomrule
    \end{tabular}}
  \end{frame}


  %--- Slide: References ---
  \begin{frame}[allowframebreaks]{Referências}
    \printbibliography[heading=none]
  \end{frame}

\end{document}
